% AppendixQuality20261023
\begingroup\tcbset{breakable/.default=false}
\chapter{电磁学补充专题}
% Editable, problem-specific diagrams for the electromagnetic appendix.
\tikzset{resistor/.style={thick,decorate,decoration={zigzag,segment length=5pt,amplitude=2pt}}}
\newcommand{\QAENamed}[3]{\begin{qualityfigure}\begingroup
  \renewcommand{\figcap}[1]{}\renewcommand{\label}[1]{}#1\endgroup
  \figcap{#2}\label{#3}\end{qualityfigure}}
\newcommand{\QAESources}{\begin{tikzpicture}[>=Latex,scale=.74,font=\small]
 \begin{scope}[xshift=-6cm]
  \draw[very thick] (0,-1.3)--(0,1.3);\node[above] at (0,1.3) {$\lambda$};
  \draw[guide] (0,0)--(1.7,0);\fill (1.7,0) circle(2pt) node[right] {$P$};
  \node[below] at (.85,0) {$a$};\node at (.8,-1.9) {无限直线};
 \end{scope}
 \begin{scope}[xshift=-2.6cm]
  \draw[very thick] (0,-.8)--(0,1.4);\coordinate(P) at(1.7,0);\coordinate(O) at(0,0);
  \draw[guide] (P)--(0,1.4) (P)--(0,-.8) (P)--(O);
  \node[left] at (0,1.4) {$T_1$};\node[left] at (0,-.8) {$T_2$};
  \node at(1.1,.3){$\theta_1$};\node at(1.05,-.25){$\theta_2$};
  \fill(P) circle(2pt) node[right] {$P$};\node[below] at(.8,0) {$a$};
  \node at (.7,-1.9) {有限直线：两端张角};
 \end{scope}
 \begin{scope}[xshift=1.4cm]
  \fill[PaleBlue] (-.8,-.4)--(.7,-.8)--(1.6,.2)--(.1,.6)--cycle;
  \draw[guide] (.4,-.1)--(.4,1.4);\node[left]at(.4,.75){$a$};\fill(.4,1.4) circle(2pt) node[right] {$P$};
  \node at(.4,-1.2) {$\sigma$};\node at(.4,-1.9) {无限平面／有限宽条带};
 \end{scope}
 \begin{scope}[xshift=6.0cm]
  \draw[fill=PaleBlue] (0,0) ellipse(.5 and 1.15);\draw[axis](0,0)--(2.1,0);
  \draw[guide](0,1.15)--(1.65,0);\fill(1.65,0) circle(2pt) node[above] {$P$};
  \node[below] at(.9,0) {$a$};\node[left] at(0,.8) {$R$};\node at(.7,-1.9) {圆盘轴线};
 \end{scope}
\end{tikzpicture}}
\newcommand{\QAEStrip}{\begin{tikzpicture}[>=Latex,scale=.85,font=\small]
 \begin{scope}[xshift=-3.8cm]
  \draw[axis](-2,0)--(2.1,0) node[right] {$x$};\draw[very thick](-1.3,0)--(1.8,0);
  \fill[PaleBlue] (.8,-.13) rectangle(1,.13);\node[below]at(.9,-.15){$\dd x$};
  \coordinate(P)at(0,1.8);\draw[guide](-1.3,0)--(P)--(1.8,0) (0,0)--(P);
  \fill(P) circle(2pt) node[above] {$P$};\node[left]at(0,.55){$a$};
  \node at(.48,1.06){$\theta_1$};\node at(-.4,1.1){$\theta_2$};
  \draw[vector](P)--(-.4,2.55) node[left]{$\dd\vect E$};
  \node[below,align=center]at(0,-.65){条带沿纸外方向无限延伸\\每条线元 $\dd\lambda=\sigma\dd x$};
 \end{scope}
 \begin{scope}[xshift=3.7cm]
  \fill[PaleBlue](-1.6,-.15)rectangle(1.6,.15);\draw[guide](-.9,-1)--(.9,-1)--(.9,1)--(-.9,1)--cycle;
  \draw[vector](0,.15)--(0,1.65)node[right]{$\vect E_+$};
  \draw[vector](0,-.15)--(0,-1.65)node[right]{$\vect E_-$};
  \node at(0,-2.15){跨平面的柱形高斯面：两个端面通量};
 \end{scope}
\end{tikzpicture}}
\newcommand{\QAEZeroField}{\begin{tikzpicture}[>=Latex,scale=.8,font=\small]
 \draw[axis](-3.8,0)--(5.3,0)node[right]{$x$};
 \fill[DeepBlue](-2,0)circle(3pt)node[below]{$+Q$};\fill[Gold](0,0)circle(3pt)node[below]{$-q$};
 \fill(3,0)circle(1.8pt)node[below right]{$x_0$};
 \draw[vector](3.05,0)--(4.5,0)node[above]{$\delta E_x>0$};\draw[vector](2.95,0)--(1.6,0);
 \draw[vector](3,1.4)--(3,.25);\draw[vector](3,-1.4)--(3,-.25);
 \draw[guide](-2,1.9)--(3,1.9) (0,1.35)--(3,1.35);
 \node[above]at(.4,1.9){$x_0+a$};\node[above]at(1.4,1.35){$x_0-a$};
 \node[below,align=center]at(0,-2){轴向扰动离开零场点，横向扰动趋向轴线\\箭头表示零场点邻域的场，非两电荷之间的全局场线};
\end{tikzpicture}}
\newcommand{\QAECavityVectors}{\begin{tikzpicture}[>=Latex,scale=.9]
 \draw[fill=PaleBlue](0,0)circle(1.8);\draw[fill=white](.65,.2)circle(.8);
 \coordinate(O)at(0,0);\coordinate(C)at(.65,.2);\coordinate(P)at(.9,.8);
 \draw[vector](O)--(C)node[midway,below]{$\vect d$};\draw[vector](C)--(P)node[right]{$\vect r'$};
 \draw[vector](O)--(P)node[midway,left]{$\vect r$};
 \foreach\p in {O,C,P}{\fill(\p)circle(1.5pt)node[below right]{$\p$};}
 \node[align=left,anchor=west]at(2.5,.4){完整大球：$+\rho$\\补偿小球：$-\rho$\\$\vect r-\vect r'=\vect d$};
\end{tikzpicture}}
\newcommand{\QAEDipoleComponents}{\begin{tikzpicture}[>=Latex,scale=.9]
 \draw[axis](-1,0)--(3.8,0)node[right]{$\vect p$};\coordinate(O)at(0,0);\coordinate(P)at(2.2,1.5);
 \draw[guide](O)--(P)node[midway,above left]{$r$};\node[below]at(O){$O$};\fill(P)circle(2pt)node[above]{$P$};
 \draw[vector](P)--++(1.1,.75)node[right]{$E_r$};\draw[vector](P)--++(-.6,.88)node[above]{$E_\theta$};
 \draw (0.65,0) arc(0:34.3:.65)node[midway,right]{$\theta$};
 \node[below,align=center]at(1,-1){极角由 $\vect p$ 向场点量起\\$E_r,E_\theta$ 的正基矢按极坐标定义};
\end{tikzpicture}}
\newcommand{\QAEPotentialPlot}{\begin{tikzpicture}[>=Latex,scale=.75,font=\small]
 \begin{scope}[xshift=-4cm]
  \draw[axis](0,-1.4)--(0,1.5)node[above]{$\varphi$};\draw[axis](0,0)--(3.2,0)node[right]{$r/r_0$};
  \draw[DeepBlue,thick,domain=.3:3,samples=60]plot(\x,{-ln(\x)});\draw[guide](1,0)--(1,-.2)node[below]{$1$};
  \node[below]at(1.6,-1.9){直线电荷：$\varphi\propto-\ln(r/r_0)$};
 \end{scope}
 \begin{scope}[xshift=3cm]
  \draw[axis](-2,0)--(2.3,0)node[right]{$z$};\draw[axis](0,-1.7)--(0,1.4)node[above]{$\varphi$};
  \draw[DeepBlue,thick](-1.6,-1.3)--(0,0)--(1.6,-1.3);
  \node[below]at(0,-1.9){面电荷：$\varphi\propto-|z|$};
 \end{scope}
\end{tikzpicture}}
\newcommand{\QAECapacitors}[1]{\begin{tikzpicture}[>=Latex,scale=.76,font=\small]
 \begin{scope}[xshift=-5cm]
  \draw[very thick](-1.1,.65)--(1.1,.65) (-1.1,-.65)--(1.1,-.65);
  \node[above]at(0,.65){$+Q,\ S=l^2$};\node[below]at(0,-.65){$-Q$};\draw[<->](1.4,-.65)--(1.4,.65)node[midway,right]{$d$};
  \ifnum#1=1 \draw[vector](0,.45)--(0,-.45)node[right]{$E$};\fi
 \end{scope}
 \begin{scope}
  \draw(0,0)circle(.7);\draw(0,0)circle(1.3);\draw[guide](0,0)--(1.3,0);
  \node[above]at(.45,0){$R_1$};\node[below]at(1,0){$R_2$};\node[left]at(-.7,0){$+Q$};\node[above]at(0,1.3){$-Q$};
  \ifnum#1=1\draw[vector](.78,-.6)--(1.03,-.8)node[right]{$E(r)$};\fi
 \end{scope}
 \begin{scope}[xshift=5cm]
  \draw(0,0)circle(.55);\draw(0,0)circle(1.25);\draw[guide](0,0)--(1.25,0);
  \node[above]at(.3,0){$a$};\node[below]at(.95,0){$b$};\node[above,align=center]at(0,1.3){同轴柱横截面\\轴向长度 $l$};
  \ifnum#1=1\draw[vector](0,-.65)--(0,-1.15)node[right]{$E(r)$};\fi
 \end{scope}
\end{tikzpicture}}
\newcommand{\QAESplitQuestion}{\begin{tikzpicture}[>=Latex,scale=.8]
 \foreach\h in {-3.5,3.5}{\draw[very thick](\h-1.6,1)--(\h+1.6,1) (\h-1.6,-1)--(\h+1.6,-1);}
 \fill[PaleBlue](-5.1,-.95)rectangle(-3.5,.95);\draw[guide](-3.5,-1)--(-3.5,1);
 \node at(-4.3,0){$\varepsilon_r$};\node at(-2.7,0){真空};\node at(-3.5,-1.5){各占面积 $S/2$};
 \fill[PaleBlue](1.9,0)rectangle(5.1,.95);\node at(3.5,.5){$\varepsilon_r$};\node at(3.5,-.5){真空};
 \draw[<->](5.4,0)--(5.4,1)node[midway,right]{$d/2$};\node at(3.5,-1.5){各层面积 $S$};
\end{tikzpicture}}
\newcommand{\QAEInsertQuestion}{\begin{tikzpicture}[>=Latex,scale=.82]
 \draw[very thick](-3.3,.7)--(3.3,.7) (-3.3,-.7)--(3.3,-.7);
 \fill[PaleBlue](-3.3,-.65)rectangle(-.3,.65);\node at(-1.8,0){$\varepsilon_r$};\node at(1.4,0){真空};
 \draw[<->](-3.3,-1.1)--(-.3,-1.1)node[midway,below]{$x$};\draw[<->](-.3,-1.1)--(3.3,-1.1)node[midway,below]{$a-x$};
 \draw[<->](3.6,-.7)--(3.6,.7)node[midway,right]{$d$};\node[above]at(0,1.1){垂直纸面宽度 $b$};
\end{tikzpicture}}
\newcommand{\QAERCQuestion}{\begin{tikzpicture}[>=Latex,scale=.8]
 \draw(-3,0)--(-2.3,0);\draw[very thick](-2.3,-.5)--(-2.3,.5);\draw(-1.95,-.3)--(-1.95,.3);
 \draw(-1.95,0)--(-.7,0);\draw[resistor](-.7,0)--(1.1,0)node[midway,above]{$R$};
 \draw(1.1,0)--(2,0);\draw[very thick](2,-.5)--(2,.5) (2.4,-.5)--(2.4,.5);
 \draw(2.4,0)--(3,0)--(3,-1.4)--(-3,-1.4)--(-3,0);
 \node[below]at(-2.15,-.55){$\mathcal E$};\node[above]at(2.2,.6){$C,\ Q(0)=0$};
\end{tikzpicture}}
\newcommand{\QAEEarthQuestion}{\begin{tikzpicture}[>=Latex,scale=.8]
 \fill[SoftGray](0,0)circle(1.4);\draw(0,0)circle(1.4);\draw[DeepBlue](0,0)circle(2);
 \draw[->](0,0)--(1.4,0)node[midway,below]{$R_\oplus$};\draw[<->](1.4,0)--(2,0)node[midway,above]{$h$};
 \node at(0,.55){地球 $-Q$};\node[above]at(0,2){等效外壳 $+Q$};
 \node[anchor=west,align=left]at(2.8,.3){薄层 $h\ll R_\oplus$\\空气电阻率 $\rho$ 均匀\\$h$ 在图中放大};
\end{tikzpicture}}
\newcommand{\QAEVacuumQuestion}{\begin{tikzpicture}[>=Latex,scale=.85]
 \draw[very thick](-3,-1)--(-3,1) (3,-1)--(3,1);
 \node[above]at(-3,1){阴极 $0$};\node[above]at(3,1){阳极 $U>0$};
 \draw[axis](-3,-1.5)--(3.8,-1.5)node[right]{$x$};\node[below]at(-3,-1.5){$0$};\node[below]at(3,-1.5){$L$};
 \draw[vector](-1,0)--(1.3,0)node[midway,above]{$e^-,\ v$};\node at(0,-2.2){电极面积 $S$；忽略初始速率与边缘效应};
\end{tikzpicture}}
\newcommand{\QAEIonQuestion}{\begin{tikzpicture}[>=Latex,scale=.85]
 \draw[very thick](-3,-1)--(-3,1) (3,-1)--(3,1);
 \node[above]at(-3,1){阳极：读出};\node[above]at(3,1){阴极：$-U_0$};
 \draw[guide](-1,-1)--(-1,1);\node[above]at(-1,1){$N$ 对};
 \draw[vector](-1,.35)--(-2.4,.35)node[midway,above]{$v^-$};\draw[vector](-1,-.35)--(1,-.35)node[midway,below]{$v^+$};
 \draw[<->](-3,-1.5)--(-1,-1.5)node[midway,below]{$x_0$};\draw[<->](-1,-1.5)--(3,-1.5)node[midway,below]{$D-x_0$};
 \draw[axis](-3,1.7)--(3.6,1.7)node[right]{$x$};
\end{tikzpicture}}
\newcommand{\QAERingQuestion}{\begin{tikzpicture}[>=Latex,scale=.85]
 \draw[thick](0,0)circle(1.3);\fill(-1.3,0)circle(2pt)node[left]{$a$};\fill(.65,1.126)circle(2pt)node[above right]{$b$};
 \draw[->](.92,.92)arc(45:20:1.3);\node[anchor=west,align=left]at(2,.3){均匀电阻丝：总电阻 $R$\\固定 $a$，沿环移动 $b$\\电源仅接于 $a,b$};
\end{tikzpicture}}
\newcommand{\QAECube}{\begin{tikzpicture}[scale=.85,font=\small]
 \coordinate(A)at(-2,-1.2);\coordinate(B)at(.8,-1.2);\coordinate(C)at(.8,1.2);\coordinate(D)at(-2,1.2);
 \coordinate(E)at(-.8,-.3);\coordinate(F)at(2,-.3);\coordinate(G)at(2,2.1);\coordinate(H)at(-.8,2.1);
 \foreach\u/\v in {A/B,B/C,C/D,D/A,E/F,F/G,G/H,H/E,A/E,B/F,C/G,D/H}{\draw[resistor](\u)--(\v);}
 \foreach\u/\pos in {A/below left,B/below,C/above left,D/left,E/below right,F/right,G/above right,H/above}{\fill(\u)circle(1.6pt)node[\pos]{$\u$};}
 \node[anchor=west,align=left]at(3,.4){每条棱均为 $R$\\三组端口：$AB,AC,AG$\\交叉处无黑点则不相连};
\end{tikzpicture}}
\newcommand{\QAELadder}[1]{\begin{tikzpicture}[scale=.85,>=Latex]
 \node[left]at(0,1){$a$};\node[left]at(0,-1){$b$};
 \foreach\i in {0,1,2}{\foreach\y in {-1,1}{\draw[resistor](2*\i,\y)--(2*\i+1.4,\y);\draw(2*\i+1.4,\y)--(2*\i+2,\y);}\draw[resistor](2*\i+1.4,-1)--(2*\i+1.4,1);}
 \draw[dotted,thick](6,1)--(6.8,1) (6,-1)--(6.8,-1);
 \ifnum#1=1\draw[dashed,DeepBlue](1.9,-1.5)rectangle(7,1.5);\node[above,DeepBlue]at(4.4,1.5){尾部仍为 $X$};\fi
 \node[below]at(3.3,-1.8){每只电阻均为 $R$；端口前没有横接电阻};
\end{tikzpicture}}
\newcommand{\QAEThreeRail}[1]{\begin{tikzpicture}[scale=.85,font=\small]
 \foreach\y/\n in {1.5/A,0/B,-1.5/C}{\node[left]at(0,\y){$\n$};}
 \foreach\i in {0,1,2}{\foreach\y in {-1.5,0,1.5}{\draw[resistor](2*\i,\y)--(2*\i+1.4,\y);\draw(2*\i+1.4,\y)--(2*\i+2,\y);}\draw[resistor](2*\i+1.4,-1.5)--(2*\i+1.4,0);\draw[resistor](2*\i+1.4,0)--(2*\i+1.4,1.5);}
 \foreach\y in {-1.5,0,1.5}{\draw[dotted,thick](6,\y)--(6.8,\y);}
 \ifnum#1=1\draw[dashed,DeepBlue](1.9,-2)rectangle(7,2);\node[above,DeepBlue]at(4.3,2){三端尾部：$p,p,q$};\fi
 \node[below]at(3,-2.3){一节含三只水平电阻和两只竖直电阻，均为 $R$};
\end{tikzpicture}}
\newcommand{\QAEThreeReduction}{\begin{tikzpicture}[>=Latex,scale=.85]
 \coordinate(A)at(-1,1.4);\coordinate(B)at(1.5,0);\coordinate(C)at(-1,-1.4);
 \draw[resistor](A)--(B)node[midway,above right]{$p/R$};\draw[resistor](B)--(C)node[midway,below right]{$p/R$};\draw[resistor](C)--(A)node[midway,left]{$q/R$};
 \foreach\n in {A,B,C}{\fill(\n)circle(2pt)node[above left]{$\n$};}
 \node[anchor=west,align=left]at(3,.4){标注量是电导\\$AB,BC$：$p/R$\\$AC$：$q/R$\\第三端口悬空，仍参与分流};
\end{tikzpicture}}
\newcommand{\QAESierpinskiPart}[1]{%
 \ifnum#1=0
  \filldraw[fill=PaleBlue,draw=DeepBlue] (0,0)--(2,0)--(1,1.7320508)--cycle;
 \else
  \pgfmathtruncatemacro{\QAENextDepth}{#1-1}
  \begin{scope}[scale=.5]\QAESierpinskiPart{\QAENextDepth}\end{scope}
  \begin{scope}[shift={(1,0)},scale=.5]\QAESierpinskiPart{\QAENextDepth}\end{scope}
  \begin{scope}[shift={(.5,.8660254)},scale=.5]\QAESierpinskiPart{\QAENextDepth}\end{scope}
 \fi}
\newcommand{\QAESierpinski}{\begin{tikzpicture}[scale=1.15,font=\small]
 \foreach\n in {0,1,2}{\begin{scope}[xshift=3.2*\n cm]
  \QAESierpinskiPart{\n}
  \node[below]at(1,-.3){第 $\n$ 阶：$3^{\n}$ 个小三角形};
 \end{scope}}
\end{tikzpicture}}



本附录补充第19—26章的电磁学专题与例题，重点讨论连续电荷分布、电势、电容器、恒定电流和复杂电阻网络。各题配有完整解答与图示，可结合正文选读。

\Needspace{18\baselineskip}

\section{静电场：连续分布与场线}

\begingroup

\graphicspath{{../../../sample/xes-original/eighteenth-section/}{../../../sample/xes-original/eighteenth-section/assets/}}

% Restrained TikZ redraws for Lesson 29.

\newcommand{\matterStructureDiagram}{%
\begin{center}
\begin{tikzpicture}[scale=0.78,>=Latex]
  % Atom
  \begin{scope}[xshift=-6.0cm]
    \fill[DeepBlue] (0,0) circle (3pt);
    \draw[TextGray,rotate=18] (0,0) ellipse (1.25 and 0.48);
    \draw[TextGray,rotate=-42] (0,0) ellipse (1.25 and 0.48);
    \draw[TextGray,rotate=78] (0,0) ellipse (1.25 and 0.48);
    \fill[Ink] (1.05,0.28) circle (2.2pt);
    \node[TextGray] at (0,-1.55) {原子 $\sim10^{-10}\,\mathrm m$};
  \end{scope}
  % Nucleus
  \begin{scope}[xshift=-2.0cm]
    \foreach \x/\y in {-0.50/0.20,-0.15/0.52,0.25/0.42,0.55/0.05,0.22/-0.35,-0.28/-0.32}{
      \fill[DeepBlue!70] (\x,\y) circle (0.28);
      \draw[Ink] (\x,\y) circle (0.28);
    }
    \node[TextGray] at (0,-1.55) {原子核 $\sim10^{-14}\,\mathrm m$};
  \end{scope}
  % Nucleon
  \begin{scope}[xshift=2.0cm]
    \draw[curve,fill=PaleBlue] (0,0) circle (0.95);
    \foreach \a in {20,80,...,320}{
      \draw[TextGray] (0,0)--({0.68*cos(\a)},{0.68*sin(\a)});
    }
    \node[TextGray] at (0,-1.55) {核子 $\sim10^{-15}\,\mathrm m$};
  \end{scope}
  % Quark content
  \begin{scope}[xshift=6.0cm]
    \draw[curve,fill=PaleBlue] (0,0) circle (0.95);
    \fill[DeepBlue] (-0.38,-0.18) circle (0.22) node[below left] {$u$};
    \fill[Gold] (0.38,-0.18) circle (0.22) node[below right] {$u$};
    \fill[Ink] (0,0.43) circle (0.22) node[above] {$d$};
    \draw[TextGray] (-0.38,-0.18)--(0.38,-0.18)--(0,0.43)--cycle;
    \node[TextGray] at (0,-1.55) {夸克尺度 $<10^{-18}\,\mathrm m$};
  \end{scope}
  \foreach \x in {-4.0,0,4.0}{\draw[vector] (\x-0.55,0)--(\x+0.55,0);}
\end{tikzpicture}

\figcap{从原子到核子内部的物质层级；长度标度为数量级示意}\label{appE:fig:s18-legacy-01}
\end{center}}

\newcommand{\coulombDiagram}{%
\begin{center}
\begin{tikzpicture}[scale=0.9,>=Latex]
  \coordinate (A) at (-2.7,0); \coordinate (B) at (2.7,0);
  \fill[DeepBlue] (A) circle (3pt) node[below=4pt] {$q_1$};
  \fill[DeepBlue] (B) circle (3pt) node[below=4pt] {$q_2$};
  \draw[guide] (A)--(B);
  \draw[<->] (-2.55,-0.75)--node[fill=white,inner sep=1pt] {$r$}(2.55,-0.75);
  \draw[vector] (A)--++(-1.6,0) node[above] {$\vect F_{21}$};
  \draw[vector] (B)--++(1.6,0) node[above] {$\vect F_{12}$};
  \node[TextGray] at (0,1.0) {图示为同号电荷；异号时两力同时反向};
\end{tikzpicture}

\figcap{点电荷间的库仑力沿两电荷连线，且满足作用力与反作用力关系}\label{appE:fig:s18-legacy-02}
\end{center}}

\newcommand{\coulombHistoryContextDiagram}{%
\begin{center}
\begin{tikzpicture}[scale=0.82,>=Latex]
  \node[inner sep=0] at (-4.4,0)
    {\includegraphics[height=3.8cm]{coulomb-portrait-redraw.png}};
  \node[TextGray] at (-4.4,-2.25) {Charles-Augustin de 库仑};
  \begin{scope}[xshift=1.2cm]
    \node[TextGray,anchor=west] at (-2.25,1.75) {俯视};
    \draw[thick] (-2.0,2.0)--(2.0,2.0);
    \draw[thick] (0,2.0)--(0,0.35);
    \draw[very thick] (-2.0,0.25)--(2.0,0.25);
    \fill[DeepBlue] (-2.0,0.25) circle (5pt) node[below] {$q_1$};
    \fill[DeepBlue] (2.0,0.25) circle (5pt) node[below] {$q_2$};
    \fill[Ink] (-2.0,1.35) circle (5pt) node[left] {$Q$};
    \draw[guide] (-2.0,1.35)--(-2.0,0.25);
    % Q is placed tangentially to the suspended arm, so the 库仑 force
    % produces a nonzero torque about the suspension axis.
    \draw[vector] (-2.0,0.25)--(-2.0,-0.62) node[left] {$\vect F$};
    \draw[<->] (0,-0.82)--node[fill=white,inner sep=1pt] {$d$}(-2.0,-0.82);
    \draw[DeepBlue,thick,->] (0,0.75) arc[start angle=90,end angle=25,radius=0.65];
    \node[DeepBlue] at (0.65,0.9) {$\theta$};
    \node[TextGray,align=left,anchor=west] at (3.0,0.55)
      {悬丝扭转\\$Fd=\kappa\theta$\\$\theta\propto r^{-2}$};
  \end{scope}
\end{tikzpicture}

\figcap{库仑与扭秤测力原理：静电力矩由悬丝的回复力矩平衡}\label{appE:fig:s18-legacy-03}
\end{center}}

\newcommand{\fieldLineDiagram}{%
\begin{center}
\begin{tikzpicture}[scale=0.72,>=Latex]
  \begin{scope}[xshift=-6.0cm]
    \fill[DeepBlue] (-0.75,0) circle (4pt) node[below] {$+Q$};
    \fill[DeepBlue] (0.75,0) circle (4pt) node[below] {$+Q$};
    \draw[->,DeepBlue,thick] (-0.98,0)--(-2.35,0);
    \draw[->,DeepBlue,thick] (0.98,0)--(2.35,0);
    \foreach \s in {-1,1}{
      \draw[->,DeepBlue,thick] ({0.75*\s},{0.18}) .. controls ({1.25*\s},0.65) and ({1.45*\s},1.45) .. ({1.55*\s},2.1);
      \draw[->,DeepBlue,thick] ({0.75*\s},{-0.18}) .. controls ({1.25*\s},-0.65) and ({1.45*\s},-1.45) .. ({1.55*\s},-2.1);
      \draw[->,DeepBlue,thick] ({0.62*\s},0.16) .. controls ({0.35*\s},0.85) and ({0.25*\s},1.55) .. ({0.18*\s},2.15);
      \draw[->,DeepBlue,thick] ({0.62*\s},-0.16) .. controls ({0.35*\s},-0.85) and ({0.25*\s},-1.55) .. ({0.18*\s},-2.15);
    }
    \node[TextGray] at (0,-2.55) {等量同号电荷};
  \end{scope}
  \begin{scope}
    \fill[DeepBlue] (-1.4,0) circle (4pt) node[below] {$+Q$};
    \fill[Ink] (1.4,0) circle (4pt) node[below] {$-Q$};
    \foreach \h in {0.28,0.62,1.0,1.45}{
      \draw[->,DeepBlue,thick]
        (-1.16,{0.55*\h}) .. controls (-0.35,{1.45*\h}) and (0.35,{1.45*\h}) .. (1.16,{0.55*\h});
      \draw[->,DeepBlue,thick]
        (-1.16,{-0.55*\h}) .. controls (-0.35,{-1.45*\h}) and (0.35,{-1.45*\h}) .. (1.16,{-0.55*\h});
    }
    \draw[->,DeepBlue,thick] (-1.05,0)--(1.05,0);
    \node[TextGray] at (0,-2.55) {等量异号电荷};
  \end{scope}
  \begin{scope}[xshift=6.0cm]
    \fill[DeepBlue] (0,0) circle (4pt) node[below right] {$q(t)$};
    \draw[guide,dashed] (0,-2.1)--(0,2.1) node[above] {振动轴};
    \draw[->,DeepBlue,thick] (0.15,0.05) .. controls (0.55,0.35) and (0.25,0.95) .. (0.95,1.30)
      .. controls (1.55,1.60) and (1.25,2.05) .. (2.25,2.25);
    \draw[->,DeepBlue,thick] (-0.15,0.05) .. controls (-0.55,0.35) and (-0.25,0.95) .. (-0.95,1.30)
      .. controls (-1.55,1.60) and (-1.25,2.05) .. (-2.25,2.25);
    \draw[->,DeepBlue,thick] (0.15,-0.05) .. controls (0.55,-0.35) and (0.25,-0.95) .. (0.95,-1.30)
      .. controls (1.55,-1.60) and (1.25,-2.05) .. (2.25,-2.25);
    \draw[->,DeepBlue,thick] (-0.15,-0.05) .. controls (-0.55,-0.35) and (-0.25,-0.95) .. (-0.95,-1.30)
      .. controls (-1.55,-1.60) and (-1.25,-2.05) .. (-2.25,-2.25);
    \node[TextGray] at (0,-2.55) {振荡电荷：场线随时间演化};
  \end{scope}
\end{tikzpicture}

\figcap{同号、异号与振荡电荷的典型电场线；第三幅不能按静电场理解}\label{appE:fig:s18-legacy-04}
\end{center}}

\newcommand{\fluxDiagram}{%
\begin{center}
\begin{tikzpicture}[scale=0.9,>=Latex]
  \coordinate (O) at (0,0);
  \coordinate (Etip) at (2.45,1.25);
  \coordinate (Stip) at (0.7,2.15);
  \draw[fill=PaleBlue,draw=DeepBlue,thick] (-1.8,-0.85)--(1.8,-0.30)--(1.8,1.05)--(-1.8,0.50)--cycle;
  \draw[vector] (O)--++(0.7,2.15) node[above] {$\vect S$};
  \draw[vector] (O)--++(2.45,1.25) node[right] {$\vect E$};
  \draw pic[draw,"$\theta$",angle radius=8mm,angle eccentricity=1.3]
    {angle=Etip--O--Stip};
  \node[TextGray] at (0,-1.35) {$\Phi_E=\vect E\cdot\vect S=ES\cos\theta$};
\end{tikzpicture}
\end{center}}

\newcommand{\gaussDiagram}{%
\begin{center}
\begin{tikzpicture}[scale=0.88,>=Latex]
  \draw[DeepBlue,thick] (0,0) ellipse (2.7 and 1.75);
  \fill[DeepBlue] (-0.45,0.15) circle (4pt) node[below right] {$q_{\rm in}$};
  \fill[Ink] (3.9,0.65) circle (4pt) node[right] {$q_{\rm out}$};
  \foreach \a in {10,55,...,325}{
    \draw[->,TextGray,thin] (-0.45,0.15)--++({1.9*cos(\a)},{1.35*sin(\a)});
  }
  \draw[vector] (2.35,0.88)--++(0.85,0.35) node[right] {$\dd\vect S$};
  \node[TextGray,align=left,anchor=west] at (5.0,0)
    {闭合面的净通量只由面内净电荷决定；\\面外电荷会改变局部电场，但净通量贡献为零。};
\end{tikzpicture}

\figcap{高斯定律中的闭合曲面}\label{appE:fig:s18-legacy-05}
\end{center}}

\newcommand{\gaussContextDiagram}{%
\begin{center}
\begin{tikzpicture}[scale=0.78,>=Latex]
  \node[inner sep=0] at (-6.0,0.2)
    {\includegraphics[width=2.2cm]{gauss-portrait-redraw.png}};
  \node[TextGray] at (-6.0,-1.65) {高斯};
  \begin{scope}[xshift=-2.8cm]
    \fill[PaleBlue] (-1.15,-0.6)--(1.15,-0.25)--(0.9,0.8)--(-1.3,0.5)--cycle;
    \draw[thick] (-1.15,-0.6)--(1.15,-0.25)--(0.9,0.8)--(-1.3,0.5)--cycle;
    \draw[vector] (0,0)--(0.55,1.55) node[above] {$\dd\vect S$};
    \draw[vector] (0,0)--(1.65,0.65) node[right] {$\vect E$};
    \coordinate (Ee) at (1.65,0.65);
    \coordinate (Oo) at (0,0);
    \coordinate (Ss) at (0.55,1.55);
    \draw pic[draw,"$\theta$",angle radius=6mm,angle eccentricity=1.3]
      {angle=Ee--Oo--Ss};
    \node[DeepBlue] at (0,-1.55) {$\dd\Phi_E=E\dd S\cos\theta$};
  \end{scope}
  \begin{scope}[xshift=1.2cm]
    \draw[DeepBlue,thick,rounded corners=12pt] plot[smooth cycle] coordinates
      {(-1.45,0.1) (-0.8,1.25) (0.7,1.1) (1.45,0.2) (0.8,-1.2) (-0.9,-1.0)};
    \fill[Ink] (0,0) circle (3pt) node[below right] {$q$};
    \foreach \a in {15,60,...,330}{
      \draw[->,TextGray] ({0.25*cos(\a)},{0.25*sin(\a)})--({1.05*cos(\a)},{0.85*sin(\a)});
    }
    \node[DeepBlue] at (0,-1.55) {$\Phi_E=q/\varepsilon_0$};
  \end{scope}
  \begin{scope}[xshift=5.35cm]
    \draw[DeepBlue,thick,rounded corners=12pt] plot[smooth cycle] coordinates
      {(-1.35,0.1) (-0.75,1.15) (0.75,1.0) (1.4,0.15) (0.75,-1.15) (-0.8,-1.0)};
    \fill[Ink] (-2.25,0) circle (3pt) node[below] {$q$};
    \foreach \y in {-0.65,-0.2,0.25,0.7}{
      \draw[->,TextGray] (-2.0,\y)--(1.0,{0.75*\y});
    }
    \node[DeepBlue] at (0,-1.55) {$\Phi_E=0$};
  \end{scope}
\end{tikzpicture}

\figcap{面积元通量、面内电荷与面外电荷三种几何情形}\label{appE:fig:s18-legacy-06}
\end{center}}

\newcommand{\gaussDifferentialCubeDiagram}{%
\begin{center}
\begin{tikzpicture}[scale=0.88,>=Latex]
  \coordinate (A) at (0,0); \coordinate (B) at (3.4,0.35);
  \coordinate (C) at (4.65,1.35); \coordinate (D) at (1.25,1.0);
  \coordinate (E) at (0,2.55); \coordinate (F) at (3.4,2.90);
  \coordinate (G) at (4.65,3.90); \coordinate (H) at (1.25,3.55);
  \draw[curve] (A)--(B)--(C)--(D)--cycle;
  \draw[curve] (E)--(F)--(G)--(H)--cycle;
  \draw[curve] (A)--(E) (B)--(F) (C)--(G) (D)--(H);
  \draw[axis] (A)--++(5.5,0.55);
  \node[below] at (5.70,0.10) {$x$};
  \draw[axis] (A)--++(1.65,1.33) node[above right] {$y$};
  \draw[axis] (A)--++(0,4.45) node[above] {$z$};
  \draw[vector] (B)--++(1.25,0.13);
  \node[anchor=west,fill=white,inner sep=0.8pt] at (5.10,1.02)
    {$E_x(x+\Delta x)$};
  \draw[TextGray,thin] (5.10,0.91)--(4.68,0.51);
  \draw[vector] (A)--++(-1.1,-0.11) node[below] {$E_x(x)$};
  \draw[vector] (H)--++(0,1.15) node[left] {$E_z(z+\Delta z)$};
  \node[TextGray,align=left,anchor=west] at (6.45,2.55)
    {六个面的净通量相加\\保留一阶差分\\$\displaystyle \nabla\!\cdot\!\vect E=\rho/\varepsilon_0$};
  \node[DeepBlue] at (2.0,-0.55) {$\Delta x$};
  \node[DeepBlue,fill=white,inner sep=0.8pt] at (4.85,1.72) {$\Delta y$};
  \node[DeepBlue] at (-0.35,1.3) {$\Delta z$};
\end{tikzpicture}

\figcap{无穷小长方体上的通量收支给出高斯定律的微分形式}\label{appE:fig:s18-legacy-07}
\end{center}}

\newcommand{\lineFieldDiagram}{%
\begin{center}
\begin{tikzpicture}[scale=0.84,>=Latex]
  \draw[very thick] (0,-2.2)--(0,2.2);
  \node[left] at (0,2.0) {$\lambda$};
  \coordinate (P) at (3.0,0);
  \fill[DeepBlue] (P) circle (2.4pt) node[right] {$P$};
  \draw[<->] (0,-0.45)--node[fill=white,inner sep=1pt] {$a$}(3.0,-0.45);
  \draw[vector] (P)--++(1.5,0) node[above] {$\vect E$};
\end{tikzpicture}
\end{center}}

\newcommand{\finiteLineDiagram}{%
\begin{center}
\begin{tikzpicture}[scale=0.9,>=Latex]
  \coordinate (O) at (0,0); \coordinate (P) at (3.2,0);
  \coordinate (T) at (0,2.3); \coordinate (D) at (0,-2.3);
  \draw[very thick] (0,-2.3)--(0,2.3);
  \fill[DeepBlue] (P) circle (2.4pt) node[right] {$P$};
  \draw[<->] (0,-0.5)--node[fill=white,inner sep=1pt] {$a$}(3.2,-0.5);
  \draw[guide] (P)--(0,2.3) (P)--(0,-2.3);
  \draw pic[draw,"$\theta_1$",angle radius=8mm,angle eccentricity=1.35]
    {angle=T--P--O};
  \draw pic[draw,"$\theta_2$",angle radius=11mm,angle eccentricity=1.25]
    {angle=O--P--D};
  \draw[vector] (P)--++(1.35,0) node[above] {$\vect E$};
\end{tikzpicture}

\figcap{有限长带电直线中垂线上的场点}\label{appE:fig:s18-legacy-08}
\end{center}}

\newcommand{\diskFieldDiagram}{%
\begin{center}
\begin{tikzpicture}[scale=0.88,>=Latex]
  \coordinate (O) at (0,0); \coordinate (P) at (4.0,0);
  \coordinate (T) at (0,2.1);
  \draw[fill=PaleBlue,draw=DeepBlue,thick] (O) ellipse (0.65 and 2.1);
  \draw[DeepBlue] (O) ellipse (0.38 and 1.2);
  \fill (O) circle (1.6pt) node[below left] {$O$};
  \fill[DeepBlue] (P) circle (2.4pt) node[right] {$P$};
  \draw[axis] (-0.8,0)--(5.3,0) node[right] {$x$};
  \draw[<->] (0,-2.5)--node[below] {$a$}(4.0,-2.5);
  \draw[guide] (P)--(0,2.1);
  \draw pic[draw,"$\theta$",angle radius=9mm,angle eccentricity=1.35]
    {angle=T--P--O};
  \node[left] at (-0.65,1.7) {$\sigma$};
\end{tikzpicture}

\figcap{均匀带电圆盘轴线上的场点；图中 $a$ 表示场点到圆盘的距离}\label{appE:fig:s18-legacy-09}
\end{center}}

\newcommand{\unequalChargeDiagram}{%
\begin{center}
\begin{tikzpicture}[scale=0.9,>=Latex]
  \coordinate (A) at (-2.1,0); \coordinate (B) at (2.1,0); \coordinate (N) at (5.4,0);
  \fill[DeepBlue] (A) circle (4pt) node[below] {$+Q$};
  \fill[Ink] (B) circle (4pt) node[below] {$-q$};
  \fill[DeepBlue] (N) circle (2.4pt) node[below] {$N$};
  \draw[axis] (-3.2,0)--(6.6,0) node[right] {$x$};
  \draw[<->] (-2.1,-0.8)--node[fill=white,inner sep=1pt] {$2a$}(2.1,-0.8);
  \draw[<->] (2.1,0.75)--node[fill=white,inner sep=1pt] {$x_0-a$}(5.4,0.75);
  \node[TextGray] at (1.65,1.45) {$Q>q>0$，零场点位于小电荷外侧};
\end{tikzpicture}

\figcap{不等量异号点电荷的零场点}\label{appE:fig:s18-legacy-10}
\end{center}}

\newcommand{\cavityDiagram}{%
\begin{center}
\begin{tikzpicture}[scale=0.9,>=Latex]
  \coordinate (O) at (0,0); \coordinate (C) at (1.15,0.40); \coordinate (P) at (1.82,1.12);
  \draw[fill=PaleBlue,draw=DeepBlue,thick] (O) circle (2.40);
  \draw[fill=white,draw=Ink,thick] (C) circle (1.02);
  \fill (O) circle (1.5pt) node[below left=2pt,fill=white,inner sep=0.5pt] {$O$};
  \fill (C) circle (1.5pt) node[below=3pt,fill=white,inner sep=0.5pt] {$C$};
  \fill[DeepBlue] (P) circle (2.2pt) node[above right=2pt,fill=white,inner sep=0.5pt] {$P$};
  \draw[vector] (O)--(C) node[pos=0.52,below=3pt,fill=white,inner sep=0.7pt] {$\vect d$};
  \draw[vector] (C)--(P) node[pos=0.55,right=4pt,fill=white,inner sep=0.7pt] {$\vect r'$};
  \draw[vector] (O)--(P) node[pos=0.48,above left=3pt,fill=white,inner sep=0.7pt] {$\vect r$};
  \node[TextGray,align=left,anchor=west] at (3.55,0.25)
    {均匀带正电球体\\叠加同位置、同大小、带负电的球体\\以挖去偏心球形空腔};
\end{tikzpicture}

\figcap{球形空腔的叠加模型，$\vect r=\vect d+\vect r'$}\label{appE:fig:s18-legacy-11}
\end{center}}

\newcommand{\multipoleDiagram}{%
\begin{center}
\begin{tikzpicture}[scale=0.9,>=Latex]
  \begin{scope}[xshift=-3.8cm]
    \coordinate (O1) at (0,0); \coordinate (X1) at (2.7,0); \coordinate (P1) at (2.7,2.3);
    \fill[Ink] (-0.75,0) circle (3pt) node[below] {$-q$};
    \fill[DeepBlue] (0.75,0) circle (3pt) node[below] {$+q$};
    \draw[<->] (-0.7,-0.75)--node[fill=white,inner sep=1pt] {$l$}(0.7,-0.75);
    \draw[axis] (0,0)--(2.7,2.3) node[right] {$P(r,\theta)$};
    \draw pic[draw,"$\theta$",angle radius=8mm,angle eccentricity=1.35]
      {angle=X1--O1--P1};
    \node[TextGray] at (0,-1.45) {电偶极子};
  \end{scope}
  \begin{scope}[xshift=4.0cm]
    \coordinate (O2) at (0,0); \coordinate (X2) at (2.7,0); \coordinate (P2) at (2.6,2.3);
    \fill[Ink] (-1.0,0) circle (3pt) node[below] {$-q$};
    \fill[DeepBlue] (0,0) circle (3pt) node[below] {$+2q$};
    \fill[Ink] (1.0,0) circle (3pt) node[below] {$-q$};
    \draw[<->] (-0.95,-0.75)--node[fill=white,inner sep=1pt] {$l$}(-0.05,-0.75);
    \draw[<->] (0.05,-0.75)--node[fill=white,inner sep=1pt] {$l$}(0.95,-0.75);
    \draw[axis] (0,0)--(2.6,2.3) node[right] {$P(r,\theta)$};
    \draw pic[draw,"$\theta$",angle radius=8mm,angle eccentricity=1.35]
      {angle=X2--O2--P2};
    \node[TextGray] at (0,-1.45) {直线电四极子};
  \end{scope}
\end{tikzpicture}

\figcap{例题 29.4 与例题 29.5 的几何关系}\label{appE:fig:s18-legacy-12}
\end{center}}

\newcommand{\dipoleDiagram}{%
\begin{center}
\begin{tikzpicture}[scale=0.95,>=Latex]
  \coordinate (O) at (0,0); \coordinate (X) at (3.3,0); \coordinate (P) at (3.3,2.5);
  \fill[Ink] (-0.85,0) circle (3pt) node[below] {$-q$};
  \fill[DeepBlue] (0.85,0) circle (3pt) node[below] {$+q$};
  \draw[<->] (-0.8,-0.75)--node[fill=white,inner sep=1pt] {$l$}(0.8,-0.75);
  \draw[axis] (O)--(P) node[right] {$P(r,\theta)$};
  \draw[guide] (O)--(X);
  \draw pic[draw,"$\theta$",angle radius=9mm,angle eccentricity=1.35]
    {angle=X--O--P};
  \draw[vector] (-0.75,0.55)--(0.75,0.55) node[midway,above] {$\vect p$};
\end{tikzpicture}

\figcap{电偶极矩 $\vect p$ 从负电荷指向正电荷}\label{appE:fig:s18-legacy-13}
\end{center}}

\newcommand{\quadrupoleDiagram}{%
\begin{center}
\begin{tikzpicture}[scale=0.95,>=Latex]
  \coordinate (O) at (0,0); \coordinate (X) at (3.3,0); \coordinate (P) at (3.1,2.5);
  \fill[Ink] (-1.1,0) circle (3pt) node[below] {$-q$};
  \fill[DeepBlue] (0,0) circle (3pt) node[below] {$+2q$};
  \fill[Ink] (1.1,0) circle (3pt) node[below] {$-q$};
  \draw[<->] (-1.05,-0.75)--node[fill=white,inner sep=1pt] {$l$}(-0.05,-0.75);
  \draw[<->] (0.05,-0.75)--node[fill=white,inner sep=1pt] {$l$}(1.05,-0.75);
  \draw[axis] (O)--(P) node[right] {$P(r,\theta)$};
  \draw[guide] (O)--(X);
  \draw pic[draw,"$\theta$",angle radius=9mm,angle eccentricity=1.35]
    {angle=X--O--P};
\end{tikzpicture}

\figcap{直线电四极子的几何关系}\label{appE:fig:s18-legacy-14}
\end{center}}

\newcommand{\slabDiagram}{%
\begin{center}
\begin{tikzpicture}[scale=0.9,>=Latex]
  \coordinate (A) at (2.65,0.52); \coordinate (B) at (3.85,0.52); \coordinate (C) at (3.85,0.35);
  \fill[PaleBlue] (-2.1,-1.15) rectangle (2.1,1.15);
  \draw[DeepBlue,thick] (-2.1,-1.15) rectangle (2.1,1.15);
  \foreach \x in {-1.7,-1.0,...,1.8}
    \foreach \y in {-0.75,0,0.75}
      \fill[DeepBlue] (\x,\y) circle (1.2pt);
  \draw[axis] (-4.0,0)--(4.0,0) node[right] {$x$};
  \node[below] at (-2.1,-1.2) {$-d$}; \node[below] at (2.1,-1.2) {$d$};
  \fill[Ink] (0,0) circle (2.8pt) node[above] {$m,q$};
  \draw[vector] (0,0)--++(1.35,0) node[above] {$\vect v_0$};
  \node[TextGray] at (0,1.55) {体电荷密度 $\rho>0$};
\end{tikzpicture}

\figcap{仅在 $-d<x<d$ 内均匀带电的无穷大厚板}\label{appE:fig:s18-legacy-15}
\end{center}}

\newcommand{\electronLensDiagram}{%
\begin{center}
\begin{tikzpicture}[scale=0.9,>=Latex]
  \draw[very thick] (0,-2.0)--(0,2.0) node[above] {正线电荷 $\lambda$};
  \draw[axis] (-4.4,0)--(4.5,0) node[right] {$x$};
  \draw[guide] (-3.8,1.05)--(3.8,1.05);
  \draw[DeepBlue,very thick,-{Latex[length=2.5mm]}]
    (-3.8,1.05) .. controls (-1.2,1.05) and (1.2,0.66) .. (3.8,0.35);
  \draw[<->] (-3.1,0)--node[left] {$b$}(-3.1,1.05);
  \draw[vector] (-3.7,1.05)--++(1.35,0) node[above] {$\vect v$};
  \draw[guide] (2.65,0.52)--++(1.2,0);
  \draw pic[draw,"$\alpha$",angle radius=8mm,angle eccentricity=1.3]
    {angle=B--A--C};
\end{tikzpicture}

\figcap{电子掠过正线电荷时受到指向金属线的微弱偏转}\label{appE:fig:s18-legacy-16}
\end{center}}


\begin{origexamplebox}{【统统算出来】}\label{appE:ex:s18-29-1}
利用积分或高斯定律计算下列均匀带电系统产生的电场：
\begin{enumerate}[label=(\arabic*),leftmargin=2.2em]
  \item 无限长直导线外距导线 $a$ 处，线电荷密度为 $\lambda$；
  \item 有限长直导线所在直线外、距垂足 $a$ 处（垂足在线段内），两端相对过垂足法线的张角分别为 $\theta_1,\theta_2$；
  \item 无限大均匀带电平面外，面电荷密度为 $\sigma$；
  \item 仅一个方向无限延伸的均匀带电条带所在平面外距平面 $a$ 处（垂足在条带内），两侧边缘所张角为
  $\theta_1,\theta_2$；
  \item 半径为 $R$、面电荷密度为 $\sigma$ 的圆盘轴线上，距圆盘中心 $a$ 处。
\end{enumerate}
\QAENamed{\QAESources}{直线、平面、条带与圆盘的题设几何；场点与电荷均不重合，距离为正。}{quality:appE:q1}
\end{origexamplebox}

\begin{xesolution}
记 $k=1/(4\pi\varepsilon_0)$。

\textbf{(1) 无限长直导线。}取与导线同轴、半径 $a$、长度 $L$ 的圆柱面。电场沿径向，
且仅圆柱侧面有通量。由高斯定律
\[
  E(2\pi aL)=\frac{\lambda L}{\varepsilon_0},
\]
故
\[
  \boxed{E=\frac{\lambda}{2\pi\varepsilon_0a}}.
\]
当 $\lambda>0$ 时方向背离导线，$\lambda<0$ 时方向指向导线。
\QAENamed{\lineFieldDiagram}{同轴高斯面只在侧面产生通量。}{quality:appE:a1-line}

\textbf{(2) 有限长直导线。}以过垂足指向场点的法线为法向正方向，沿导线取坐标 $x$，电荷元
$\dd q=\lambda\dd x$ 到场点的距离为 $r=\sqrt{a^2+x^2}$。它产生的场元为
\[
  \dd E=\frac{k\lambda\dd x}{a^2+x^2}.
\]
将场元分解到法向和导线方向，并令 $x=a\tan\theta$，得到
\[
  \boxed{E_\perp=\frac{k\lambda}{a}
  (\sin\theta_1+\sin\theta_2)},
\]
\[
  \boxed{E_\parallel=\frac{k\lambda}{a}
  (\cos\theta_1-\cos\theta_2)}.
\]
$E_\parallel$ 的正方向取向张角为 $\theta_1$ 的一端；若线段关于过垂足的法线对称，
$\theta_1=\theta_2=\theta$，则切向分量抵消，
\[
  \boxed{E=\frac{2k\lambda}{a}\sin\theta}.
\]
令 $\theta\to\pi/2$，即回到无限长直导线的结果。

\textbf{(3) 无限大平面。}取跨过平面的薄圆柱高斯面。由平面对称性，两底面的电场大小
相同且垂直平面，侧面无通量，因而
\[
  2ES=\frac{\sigma S}{\varepsilon_0},
  \qquad
  \boxed{E=\frac{\sigma}{2\varepsilon_0}}.
\]

\textbf{(4) 无限长条带。}把条带分成许多与无限方向平行的细线，每条细线的线电荷密度
$\dd\lambda=\sigma\dd x$。设场点到条带所在平面的垂足为坐标原点，条带两边分别位于
$x=-a\tan\theta_2$ 和 $x=a\tan\theta_1$。叠加无限长直线的电场，有
\[
  \boxed{E_\perp=\frac{\sigma}{2\pi\varepsilon_0}
  (\theta_1+\theta_2)},
\]
\[
  \boxed{E_\parallel=
  \frac{\sigma}{2\pi\varepsilon_0}
  \ln\frac{\cos\theta_1}{\cos\theta_2}}.
\]
这里 $E_\parallel$ 的正方向取向张角为 $\theta_1$ 的一侧。
若场点位于条带中垂面，$\theta_1=\theta_2=\theta$，则
\[
  \boxed{E=\frac{\sigma\theta}{\pi\varepsilon_0}}.
\]
当条带宽度趋于无穷，$\theta\to\pi/2$，再次得到 $E=\sigma/(2\varepsilon_0)$。

\textbf{(5) 均匀带电圆盘。}把圆盘分成半径 $s$、宽度 $\dd s$ 的同心圆环，
$\dd q=2\pi\sigma s\dd s$。圆环在轴线场点处的横向分量相消，轴向场元为
\[
  \dd E_x=k\frac{a\,\dd q}{(a^2+s^2)^{3/2}}.
\]
从 $s=0$ 积分到 $R$，得到
\[
  \boxed{E_x=\frac{\sigma}{2\varepsilon_0}
  \left(1-\frac{a}{\sqrt{a^2+R^2}}\right)}.
\]
若圆盘边缘对场点的半张角为 $\theta$，则
\[
  \boxed{E_x=\frac{\sigma}{2\varepsilon_0}(1-\cos\theta)}.
\]
\QAENamed{\diskFieldDiagram}{圆盘分割为同心圆环，各环横向分量相消。}{quality:appE:a1-disk}

\QAENamed{\QAEStrip}{条带的线元分割与无限平面的双端面高斯面。}{quality:appE:a1}
图~\ref{quality:appE:a1}左图用 $x=a\tan\theta$ 将条带线元换成角度；每条线元的水平场沿场点到线元方向的反向，因此积分后水平分量含对数而非角度和。右图两个端面都要计入通量。图~\ref{quality:appE:a1-line}与图~\ref{quality:appE:a1-disk}分别对应同轴高斯法和同心环积分法。
\end{xesolution}

\begin{origexamplebox}{【奇妙的电场线】}\label{appE:ex:s18-29-2}
点电荷 $+Q$、$-q$ 分别位于 $(-a,0,0)$、$(a,0,0)$，且 $Q>q>0$。系统关于两电荷连线
旋转对称。回答：
\begin{enumerate}[label=(\arabic*),leftmargin=2.2em]
  \item $+Q$ 发出的电场线能否到达无穷远？$-q$ 接收的电场线是否来自无穷远？
  \item 空间中是否存在电场为零的点？该点附近的电场线有何特征？
  \item 无穷远处的电场有什么特征？
  \item 写出电场线方程的一种隐式表示。
\end{enumerate}
\QAENamed{\unequalChargeDiagram}{不等量异号电荷位于同一轴线上，$Q>q>0$；场线数仅表示通量比例。}{quality:appE:q2}
\end{origexamplebox}

\begin{xesolution}
\textbf{(1)} 电场线条数只能用来表示通量比例，不是绝对数量。由于净电荷为 $Q-q>0$，
$+Q$ 发出的电通量中有一部分终止于 $-q$，其余必然延伸到无穷远。$-q$ 所接收的全部电场线
都可来自 $+Q$，不需要由无穷远进入。

\textbf{(2)} 零场点只能在两电荷连线延长线上。两电荷之间的电场同向，不可能抵消；
在 $+Q$ 左侧，大电荷又更近，也不能抵消。因此零场点在 $-q$ 的右侧。设其坐标为 $x_0>a$，
则
\[
  \frac{kQ}{(x_0+a)^2}=\frac{kq}{(x_0-a)^2},
\]
解得
\[
  \boxed{x_0=a\frac{\sqrt Q+\sqrt q}{\sqrt Q-\sqrt q}}.
\]
零场点附近的场具有鞍点结构：沿两电荷连线方向，电场线从零场点两侧离开；沿垂直方向，
电场线向该轴线汇聚。电场线本身不穿过电场为零的点，因为该点没有确定的切向方向。

\textbf{(3)} 当 $r\gg a$ 时，观察者分辨不出两电荷的间距，最低阶电场等效于位于原点的
净电荷 $Q-q$：
\[
  \boxed{\vect E\simeq k\frac{Q-q}{r^2}\uvec r}.
\]
因此无穷远处的电场线近似沿径向向外，偶极修正比主项多衰减一个 $1/r$。

\textbf{(4)} 在包含两电荷的任一子午面内，设场点相对于 $+Q$、$-q$ 的极角分别为
$\theta_+$、$\theta_-$，二者的极轴都取沿 $+x$ 方向。由轴对称电场的通量函数可得，
每一条电场线满足
\[
  \boxed{Q(1-\cos\theta_+)-q(1-\cos\theta_-)=C},
\]
其中 $C$ 由所选电场线决定。等价地，也可写成
\[
  Q\cos\theta_+-q\cos\theta_-=C'.
\]
这是一种隐式方程；绕 $x$ 轴旋转即可得到三维电场线族。

\QAENamed{\QAEZeroField}{零场点两侧的轴向场与垂直扰动方向；只画零场点邻域的线性结构。}{quality:appE:a2}

图~\ref{quality:appE:a2}中较小负电荷距 $x_0$ 更近，才能与较大的正电荷抵消。令 $A=\partial E_x/\partial x|_{x_0}>0$，无源邻域的轴对称性和 $\nabla\cdot\vect E=0$ 给出横向线性系数 $-A/2$，从而得到图示鞍形方向。电荷之间并不存在这样的零场点。
\end{xesolution}

\begin{origexamplebox}{【球壳空洞】}\label{appE:ex:s18-29-3}
一个球体内部挖去偏心球形空腔，未挖去部分具有均匀体电荷密度 $\rho$。
\begin{enumerate}[label=(\arabic*),leftmargin=2.2em]
  \item 求空腔内的电场；
  \item 分析它与万有引力中同类模型的联系，并讨论题目若把球体称为“金属球体”会有什么问题。
\end{enumerate}
\QAENamed{\cavityDiagram}{均匀带电绝缘球中的偏心球形空腔；空腔完全处于大球内。}{quality:appE:q3}
\end{origexamplebox}

\begin{xesolution}
设完整大球球心为 $O$，空腔球心为 $C$，$\vect d=\overrightarrow{OC}$。把实际电荷分布分解为：
以 $O$ 为球心、密度为 $\rho$ 的完整均匀带电球，加上以 $C$ 为球心、密度为 $-\rho$ 的小球。
均匀带电实心球内部距球心位置矢量为 $\vect r$ 的点，电场为
\[
  \vect E=\frac{\rho}{3\varepsilon_0}\vect r.
\]
对空腔内任一点 $P$，记 $\vect r=\overrightarrow{OP}$、$\vect r'=\overrightarrow{CP}$，
则 $\vect r=\vect d+\vect r'$。叠加两球的电场：
\[
  \vect E=\frac{\rho}{3\varepsilon_0}\vect r
  -\frac{\rho}{3\varepsilon_0}\vect r'
  =\boxed{\frac{\rho}{3\varepsilon_0}\vect d}.
\]
所以空腔内电场是匀强电场，当 $\rho>0$ 时方向由大球球心指向空腔球心。

万有引力也服从平方反比定律，均匀球形空腔可用完全相同的叠加方法计算，只需把电学常量与
质量密度替换为引力常量与质量密度；引力只有吸引，最终场方向相应改变。

若把球体称为“金属球体”并同时规定内部均匀体电荷密度，这在静电平衡下不成立。导体内
自由电荷会重新分布，导体材料内部电场必须为零。因此上述答案适用于均匀体电荷被冻结的
绝缘球体；若确为处于静电平衡的导体，必须按导体边界条件重新求解。

\QAENamed{\QAECavityVectors}{带正、负密度的两个完整球叠加；$\vect r=\vect d+\vect r'$。}{quality:appE:a3}

图~\ref{quality:appE:a3}的两条位置矢量都终止于同一 $P$，相减后只剩与 $P$ 无关的 $\vect d$，这正是空腔内场均匀的原因。若 $\rho<0$，方向随密度符号反转；万有引力对应 $\vect g=-4\pi G\rho_m\vect d/3$。
\end{xesolution}

\begin{origexamplebox}{【电偶极子】}\label{appE:ex:s18-29-4}
一正一负两个点电荷的电荷量分别为 $+q$、$-q$，间距为 $l$。求 $r\gg l$ 处电场强度的
大小和方向，并求在该近似下的电场线方程。
\QAENamed{\dipoleDiagram}{两个等量异号点电荷与远场位置；偶极矩方向由负电荷指向正电荷。}{quality:appE:q4}
\end{origexamplebox}

\begin{xesolution}
定义电偶极矩
\[
  \vect p=q\vect l,
\]
方向由负电荷指向正电荷。取偶极子中心为极点，以 $\vect p$ 方向为极轴。当 $r\gg l$ 时，
两点电荷势的最低阶非零项为
\[
  \varphi(r,\theta)=k\frac{p\cos\theta}{r^2}.
\]
由 $\vect E=-\nabla\varphi$ 得
\[
  \boxed{E_r=\frac{2kp\cos\theta}{r^3}},
  \qquad
  \boxed{E_\theta=\frac{kp\sin\theta}{r^3}}.
\]
因此
\[
  \boxed{E=\frac{kp}{r^3}\sqrt{1+3\cos^2\theta}}.
\]
方向由 $E_r,E_\theta$ 共同确定；轴线上电场沿 $\vect p$，中垂面上电场与 $\vect p$ 反向。

电场线满足
\[
  \frac{\dd r}{r\dd\theta}=\frac{E_r}{E_\theta}=2\cot\theta.
\]
积分后得到
\[
  \boxed{r=C\sin^2\theta},
\]
其中 $C>0$ 标记不同的电场线。

\QAENamed{\QAEDipoleComponents}{远场中的径向、极角方向及电场分量；图示第一象限。}{quality:appE:a4}

图~\ref{quality:appE:a4}中 $\theta$ 沿从正极轴转向 $P$ 的方向增大，故 $E_\theta$ 的正方向是逆时针的局部切向。把 $E_r,E_\theta$ 按这两条正交基矢合成，才能得到空间方向。$r=C\sin^2\theta$ 只在理想偶极或远场有效，曲线向 $r\sim l$ 延伸的部分不再代表实际两电荷的近场。
\end{xesolution}

\endgroup

\Needspace{18\baselineskip}

\section{电势与相互作用能}

\begingroup

\graphicspath{{../../../sample/xes-original/nineteenth-section/}{../../../sample/xes-original/nineteenth-section/assets/}}

% Restrained TikZ redraws for Lesson 30.

\newcommand{\potentialGradientDiagram}{%
\begin{center}
\begin{tikzpicture}[scale=0.72,>=Latex]
  \begin{scope}[xshift=-6.0cm]
    \foreach \x in {-1.8,-0.9,0,0.9,1.8}{\draw[DeepBlue,thick] (\x,-1.65)--(\x,1.65);}
    \foreach \y in {-1.15,-0.38,0.38,1.15}{\draw[->,TextGray] (-2.15,\y)--(2.15,\y);}
    \node[TextGray] at (0,-2.1) {匀强电场};
  \end{scope}
  \begin{scope}
    \fill[DeepBlue] (0,0) circle (4pt) node[below right] {$+Q$};
    \foreach \r in {0.55,1.05,1.55}{\draw[DeepBlue,thick] (0,0) circle (\r);}
    \foreach \a in {0,45,...,315}{\draw[->,TextGray] ({0.35*cos(\a)},{0.35*sin(\a)})--({2.05*cos(\a)},{2.05*sin(\a)});}
    \node[TextGray] at (0,-2.1) {孤立正点电荷};
  \end{scope}
  \begin{scope}[xshift=6.0cm]
    \fill[DeepBlue] (-1.15,0) circle (4pt) node[below] {$+q$};
    \fill[Ink] (1.15,0) circle (4pt) node[below] {$-q$};
    \foreach \h in {0.35,0.75,1.15}{
      \draw[->,TextGray] (-0.95,{0.55*\h}) .. controls (-0.25,{1.45*\h}) and (0.25,{1.45*\h}) .. (0.95,{0.55*\h});
      \draw[->,TextGray] (-0.95,{-0.55*\h}) .. controls (-0.25,{-1.45*\h}) and (0.25,{-1.45*\h}) .. (0.95,{-0.55*\h});
    }
    \draw[DeepBlue,thick,dashed] (0,-1.85)--(0,1.85) node[above] {$\varphi=0$};
    \node[TextGray] at (0,-2.1) {等量异号点电荷};
  \end{scope}
\end{tikzpicture}

\figcap{三类典型电场的等势线（蓝）与电场线（灰）；两者处处正交}\label{appE:fig:s19-legacy-01}
\end{center}}

\newcommand{\dipoleEquipotentialDiagram}{%
\begin{center}
\begin{tikzpicture}[scale=0.82,>=Latex]
  \fill[DeepBlue] (-1.35,0) circle (4pt) node[below] {$+q$};
  \fill[Ink] (1.35,0) circle (4pt) node[below] {$-q$};
  \foreach \h in {0.35,0.7,1.05,1.4}{
    \draw[DeepBlue,thick] (-1.12,{0.5*\h}) .. controls (-0.35,{1.45*\h}) and (0.35,{1.45*\h}) .. (1.12,{0.5*\h});
    \draw[DeepBlue,thick] (-1.12,{-0.5*\h}) .. controls (-0.35,{-1.45*\h}) and (0.35,{-1.45*\h}) .. (1.12,{-0.5*\h});
  }
  \draw[TextGray,thick,dashed] (0,-2.4)--(0,2.4) node[above] {$\varphi=0$};
  \draw[TextGray] (-1.35,0) circle (0.55);
  \draw[TextGray] (1.35,0) circle (0.55);
\end{tikzpicture}

\figcap{等量异号点电荷的电场线与零等势面}\label{appE:fig:s19-legacy-02}
\end{center}}

\newcommand{\ionizationEnergyDiagram}{%
\begin{center}
\begin{tikzpicture}[x=0.54cm,y=0.12cm,>=Latex]
  \draw[axis] (0,0)--(19.4,0) node[right] {原子序数 $Z$};
  \draw[axis] (0,0)--(0,28) node[above] {第一电离能/eV};
  \foreach \y in {5,10,15,20,25}{
    \draw[TextGray,thin] (0,\y)--(18.6,\y);
    \node[TextGray,left] at (0,\y) {\y};
  }
  \draw[DeepBlue,thick] plot coordinates {
    (1,13.60) (2,24.59) (3,5.39) (4,9.32) (5,8.30) (6,11.26)
    (7,14.53) (8,13.62) (9,17.42) (10,21.56) (11,5.14) (12,7.65)
    (13,5.99) (14,8.15) (15,10.49) (16,10.36) (17,12.97) (18,15.76)};
  \foreach \x/\y in {1/13.60,2/24.59,3/5.39,4/9.32,5/8.30,6/11.26,7/14.53,8/13.62,9/17.42,10/21.56,11/5.14,12/7.65,13/5.99,14/8.15,15/10.49,16/10.36,17/12.97,18/15.76}{
    \fill[DeepBlue] (\x,\y) circle (2pt);
  }
  \foreach \x/\lab in {1/H,2/He,3/Li,6/C,8/O,10/Ne,11/Na,14/Si,17/Cl,18/Ar}{
    \node[TextGray,below,rotate=55,anchor=east] at (\x,-0.3) {\lab};
  }
  \node[TextGray,align=left,anchor=west] at (19.9,16)
    {稀有气体形成峰值\\碱金属形成低谷\\多数原子处于 $5$--$25\,\mathrm{eV}$};
\end{tikzpicture}

\figcap{前18号元素的第一电离能；折线连接仅用于显示周期性趋势}\label{appE:fig:s19-legacy-03}
\end{center}}

\newcommand{\sphereEnergyDiagram}{%
\begin{center}
\begin{tikzpicture}[scale=0.88,>=Latex]
  \coordinate (O) at (0,0);
  \draw[curve] (O) circle (1.55);
  \foreach \a in {0,30,...,330}{
    \fill[DeepBlue] ({1.55*cos(\a)},{1.55*sin(\a)}) circle (1.2pt);
  }
  \draw[vector] (O)--(1.1,1.1) node[midway,above left] {$R$};
  \draw[vector] (1.55,0)--++(1.4,0) node[above] {$E(R)$};
  \draw[vector] (1.05,1.14)--++(0.95,1.02) node[right] {$p$};
  \node[TextGray,align=left,anchor=west] at (4.1,0.35)
    {$Q$ 均匀分布在球面\\表面电场 $E(R)=Q/(4\pi\varepsilon_0R^2)$\\静电压强 $p=\varepsilon_0E^2/2$};
\end{tikzpicture}

\figcap{均匀带电球壳的电场能量与静电压强}\label{appE:fig:s19-legacy-04}
\end{center}}

\newcommand{\zeroPotentialSphereDiagram}{%
\begin{center}
\begin{tikzpicture}[scale=0.85,>=Latex]
  \draw[axis] (-4.6,0)--(5.25,0) node[right] {$x$};
  \fill[Ink] (0,0) circle (4pt) node[below] {$-q$};
  \fill[DeepBlue] (4.2,0) circle (4pt) node[below] {$+Q$};
  % Draw the exact Q=2q member of the Apollonius family: centre=-a/3,
  % radius=2a/3.  It therefore encloses the smaller charge at the origin.
  \draw[DeepBlue,thick] (-1.4,0) circle (2.8);
  \node[TextGray] at (-1.4,3.25) {$\varphi=0$（示意取 $Q=2q$）};
  \draw[<->] (0,-0.8)--node[fill=white,inner sep=1pt] {$a$}(4.2,-0.8);
  \node[TextGray] at (-1.4,-3.25)
    {当 $Q>q$ 时，零等势面包围较小的负电荷};
\end{tikzpicture}

\figcap{不等量异号点电荷的零等势面是 Apollonius 球面}\label{appE:fig:s19-legacy-05}
\end{center}}

\newcommand{\twoLineDiagram}{%
\begin{center}
\begin{tikzpicture}[scale=0.9,>=Latex]
  \draw[axis] (-4.6,0)--(4.6,0) node[right] {$x$};
  \draw[axis] (0,-2.7)--(0,2.7) node[above] {$y$};
  \fill[DeepBlue] (1.8,0) circle (4pt) node[below] {$+\lambda$};
  \fill[Ink] (-1.8,0) circle (4pt) node[below] {$-\lambda$};
  \draw[TextGray,dashed] (0,-2.4)--(0,2.4) node[above right] {$\varphi=0$};
  % Equipotentials: Apollonius circles surrounding one line charge.
  % K=3 and K=5 members.  Every circle surrounds the corresponding
  % line charge at x=+/-a rather than sitting between the charges.
  \draw[DeepBlue,thick] (2.25,0) circle (1.35);
  \draw[DeepBlue,thick] (1.95,0) circle (0.75);
  \draw[DeepBlue,thick] (-2.25,0) circle (1.35);
  \draw[DeepBlue,thick] (-1.95,0) circle (0.75);
  % Field lines: circles through both line charges (only representative arcs).
  \draw[TextGray,thick,dashed] (0,1.65) circle (2.44);
  \draw[TextGray,thick,dashed] (0,-1.65) circle (2.44);
  \node[TextGray] at (0,-3.2) {实线是等势线；虚线是与它们正交、并通过两导线的电场线};
\end{tikzpicture}

\figcap{两根异号无限长带电直线的横截面}\label{appE:fig:s19-legacy-06}
\end{center}}

\newcommand{\finiteLineDiagramThirty}{%
\begin{center}
\begin{tikzpicture}[scale=0.92,>=Latex]
  \draw[axis] (-4.2,0)--(4.2,0) node[right] {$x$};
  \draw[axis] (0,-2.4)--(0,2.4) node[above] {$\rho$};
  \draw[DeepBlue,very thick] (-1.7,0)--(1.7,0);
  \fill (0,0) circle (1.5pt) node[below] {$O$};
  \node[below] at (-1.7,0) {$-a$}; \node[below] at (1.7,0) {$a$};
  \coordinate (P) at (2.7,1.6);
  \fill[DeepBlue] (P) circle (2.5pt) node[right] {$P(x,\rho)$};
  \draw[guide] (P)--(-1.7,0) (P)--(1.7,0);
  \node[TextGray] at (0,-2.95) {$r_++r_-=\text{常量}$ 为旋转椭球；$r_--r_+=\text{常量}$ 为旋转双曲面};
\end{tikzpicture}

\figcap{有限长均匀带电直线的轴对称几何}\label{appE:fig:s19-legacy-07}
\end{center}}

\newcommand{\repulsionDiagram}{%
\begin{center}
\begin{tikzpicture}[scale=0.95,>=Latex]
  \draw[thick] (-1.8,3.5)--(1.8,3.5);
  \foreach \x in {-1.6,-1.2,...,1.6}{\draw[TextGray] (\x,3.5)--++(-0.18,0.18);}
  \coordinate (O) at (0,3.5); \coordinate (A) at (-1.55,1.05); \coordinate (B) at (0,0.42);
  \draw[curve] (O)--(A);
  \fill[DeepBlue] (A) circle (4pt) node[below left=3pt,fill=white,inner sep=0.7pt] {$A(+q)$};
  \fill[DeepBlue] (B) circle (4pt) node[right=4pt,fill=white,inner sep=0.7pt] {$B(+q)$};
  \draw[guide] (O)--(B);
  \draw[<->] (-2.25,0.42)--node[left] {$h$}(-2.25,1.05);
  \draw[vector] (A)--++(0,-1.15) node[below right=2pt,fill=white,inner sep=0.7pt] {$m\vect g$};
  \draw[vector] (A)--++(0.75,1.18) node[above right=2pt,fill=white,inner sep=0.7pt] {$\vect T$};
  \draw[vector] (A)--++(-1.25,0.51) node[above left=2pt,fill=white,inner sep=0.7pt] {$\vect F_e$};
  \node[TextGray,align=left,anchor=west] at (3.0,1.55)
    {$B$ 最终位于 $A$ 原来的最低点\\慢速移动时可用能量关系};
\end{tikzpicture}

\figcap{同号小球排斥过程的始末几何}\label{appE:fig:s19-legacy-08}
\end{center}}

\newcommand{\bohrDiagram}{%
\begin{center}
  \begin{tikzpicture}[scale=0.9,>=Latex]
  \fill[DeepBlue] (0,0) circle (5pt) node[below left=4pt,fill=white,inner sep=0.7pt] {$p(+e)$};
  \draw[DeepBlue,thick] (0,0) circle (2.05);
  \fill[Ink] (2.05,0) circle (4pt) node[below right=4pt,fill=white,inner sep=0.7pt] {$e^-(-e)$};
  \draw[vector] (2.05,0)--++(0,1.45) node[above right=2pt,fill=white,inner sep=0.7pt] {$\vect v$};
  \draw[vector] (2.05,0)--++(-1.25,0) node[pos=0.55,above=4pt,fill=white,inner sep=0.7pt] {$\vect F_e$};
  \draw[<->] (0,-1.05)--node[fill=white,inner sep=1pt] {$r$}(2.05,-1.05);
\end{tikzpicture}

\figcap{玻尔基态模型中的电子圆周运动}\label{appE:fig:s19-legacy-09}
\end{center}}

\newcommand{\rutherfordDiagram}{%
\begin{center}
\begin{tikzpicture}[scale=0.9,>=Latex]
  \fill[DeepBlue] (0,0) circle (5pt) node[below right] {$Ze$};
  \draw[axis] (-4.8,0)--(4.8,0) node[right] {$x$};
  \draw[guide] (-4.2,1.25)--(0,1.25);
  \draw[DeepBlue,very thick,-{Latex[length=2.7mm]}]
    (-4.2,1.25) .. controls (-1.5,1.25) and (-0.8,2.2) .. (3.4,3.0);
  \draw[<->] (-3.4,0)--node[left] {$b$}(-3.4,1.25);
  \draw[vector] (-4.0,1.25)--++(1.3,0) node[above] {$\vect v_\infty$};
  \draw[guide] (1.9,2.55)--++(1.7,0);
  \coordinate (A) at (1.9,2.55); \coordinate (B) at (3.6,2.55); \coordinate (C) at (3.4,3.0);
  \draw pic[draw,"$\theta$",angle radius=9mm,angle eccentricity=1.25]
    {angle=B--A--C};
\end{tikzpicture}

\figcap{库仑斥力下的双曲线散射轨迹}\label{appE:fig:s19-legacy-10}
\end{center}}


\begin{origexamplebox}{【电势算电场】}\label{appE:ex:s19-30-1}
利用“先求电势，再求电场”的方法，重新计算上一讲中的典型带电系统：点电荷系、有限带电体
和无限长或无限大的电荷分布。分析哪些情况下不能取无穷远为零电势，并给出处理方法。

\QAENamed{\QAESources}{有限带电体与无限分布的比较；先确定几何和有限参考点，再求电势差。}{quality:appE:q5}
\end{origexamplebox}

\begin{xesolution}
方法的核心是
\[
  \varphi(\vect r)=k\int\frac{\dd q}{|\vect r-\vect r'|},
  \qquad
  \vect E=-\nabla\varphi,
  \qquad k=\frac1{4\pi\varepsilon_0}.
\]

\textbf{点电荷系。}直接标量叠加
\[
  \varphi=k\sum_i\frac{Q_i}{r_i},
\]
然后求梯度即得
\[
  \vect E=k\sum_i\frac{Q_i}{r_i^3}\vect r_i.
\]
这种方法对偶极子、多极子尤其方便，因为远场展开先在标量电势上进行，再求导即可。

\textbf{均匀带电圆盘。}半径 $R$、面密度 $\sigma$ 的圆盘，在轴线上距圆心 $z>0$ 处的电势为
\[
  \varphi(z)=k\int_0^R\frac{2\pi\sigma s\dd s}{\sqrt{z^2+s^2}}
  =\frac{\sigma}{2\varepsilon_0}
  \bigl(\sqrt{z^2+R^2}-z\bigr).
\]
因此
\[
  \boxed{E_z=-\frac{\dd\varphi}{\dd z}
  =\frac{\sigma}{2\varepsilon_0}
  \left(1-\frac{z}{\sqrt{z^2+R^2}}\right)}.
\]

\textbf{无限长直线。}若线电荷密度为 $\lambda$，$E_r=\lambda/(2\pi\varepsilon_0r)$。
从有限参考半径 $r_0$ 积分得
\[
  \boxed{\varphi(r)-\varphi(r_0)
  =-\frac{\lambda}{2\pi\varepsilon_0}\ln\frac r{r_0}}.
\]
因为 $r\to\infty$ 时对数发散，所以不能令无穷远处电势为零。

\textbf{无限大平面。}若面电荷密度为 $\sigma$，取平面为 $z=0$，则
\[
  E_z=\frac{\sigma}{2\varepsilon_0}\operatorname{sgn}(z),
\]
取 $\varphi(0)=0$ 后
\[
  \boxed{\varphi(z)=-\frac{\sigma}{2\varepsilon_0}|z|}.
\]
它同样不能用无穷远作为零电势。对非局域电荷分布，正确做法是选定有限参考点，始终只计算
电势差；若多个无限分布组合后总场在远处衰减足够快，也可先把组合视为整体再选零点。

\QAENamed{\QAEPotentialPlot}{正电荷密度时的电势形状；两图的纵向正尺度分别为 $\lambda/(2\pi\varepsilon_0)$、$\sigma/(2\varepsilon_0)$。}{quality:appE:a5}

图~\ref{quality:appE:a5}中对数和绝对值曲线在远处均不趋于有限常数，所以不能给它们硬加“无穷远零电势”。求导会消去有限参考点带来的常数，电场不随零点选择而变。平面 $z=0$ 处电势连续、法向导数跳变，与面电荷边界相符。
\end{xesolution}

\endgroup

\Needspace{18\baselineskip}

\section{电介质与电容器}

\begingroup

\graphicspath{{../../../sample/xes-original/twenty-first-section/}{../../../sample/xes-original/twenty-first-section/assets/}}

% Restrained TikZ redraws for Lesson 32.

\newcommand{\arbitraryCapacitorDiagram}{%
\begin{center}
\begin{tikzpicture}[x=1cm,y=1cm,>=Latex]
  \fill[PaleBlue] plot[smooth cycle,tension=0.78] coordinates
    {(-3.7,0.1) (-3.3,1.35) (-2.05,1.55) (-1.45,0.55) (-1.8,-0.75) (-3.05,-1.25) (-3.95,-0.55)};
  \draw[DeepBlue,very thick] plot[smooth cycle,tension=0.78] coordinates
    {(-3.7,0.1) (-3.3,1.35) (-2.05,1.55) (-1.45,0.55) (-1.8,-0.75) (-3.05,-1.25) (-3.95,-0.55)};
  \fill[SoftGray] plot[smooth cycle,tension=0.78] coordinates
    {(3.7,0.15) (3.15,1.45) (1.85,1.25) (1.4,0.15) (1.95,-1.05) (3.35,-1.25) (4.0,-0.45)};
  \draw[Ink,very thick] plot[smooth cycle,tension=0.78] coordinates
    {(3.7,0.15) (3.15,1.45) (1.85,1.25) (1.4,0.15) (1.95,-1.05) (3.35,-1.25) (4.0,-0.45)};
  \foreach \p in {(-3.35,0.65),(-2.65,1.05),(-2.15,0.15),(-3.1,-0.55)}
    \node[DeepBlue] at \p {$+$};
  \foreach \p in {(3.3,0.7),(2.45,0.95),(1.95,0.05),(2.85,-0.55)}
    \node[Ink] at \p {$-$};
  \draw[component] (-1.65,0.82) .. controls (-0.5,1.5) and (0.5,1.5) .. (1.65,0.82);
  \draw[component] (-1.45,0.15) .. controls (-0.45,0.65) and (0.45,0.65) .. (1.45,0.15);
  \draw[component] (-1.62,-0.55) .. controls (-0.5,-1.15) and (0.5,-1.15) .. (1.72,-0.55);
  \node[DeepBlue] at (-2.7,-1.65) {$+Q,\ \varphi_+$};
  \node[Ink] at (2.7,-1.65) {$-Q,\ \varphi_-$};
  \draw[<->,TextGray] (-1.05,-1.9)--node[below] {$U=\varphi_+-\varphi_-$}(1.05,-1.9);
\end{tikzpicture}

\figcap{任意形状的两块导体组成双极电容器}\label{appE:fig:s21-legacy-01}
\end{center}}

\newcommand{\capacitorGeometryDiagram}{%
\begin{center}
\begin{tikzpicture}[scale=0.82,>=Latex]
  % parallel plates
  \draw[DeepBlue,very thick] (-5.2,0.8)--(-2.6,0.8);
  \draw[DeepBlue,very thick] (-5.2,-0.8)--(-2.6,-0.8);
  \foreach \x in {-4.9,-4.4,...,-2.9}{
    \node[DeepBlue] at (\x,0.48) {$+$}; \node[Ink] at (\x,-0.48) {$-$};
    \draw[component] (\x,0.32)--(\x,-0.32);
  }
  \draw[<->] (-5.55,-0.8)--node[left] {$d$}(-5.55,0.8);
  \node[TextGray] at (-3.9,-1.45) {平行板};
  % spherical
  \draw[curve] (0,0) circle (0.73); \draw[curve] (0,0) circle (1.28);
  \draw[<->] (0,0)--node[above] {$R_1$}(0.73,0);
  \draw[<->] (0,-0.28)--node[below] {$R_2$}(1.28,-0.28);
  \node[TextGray] at (0,-1.45) {同心球};
  % coaxial
  \draw[curve] (3.6,0) circle (0.62); \draw[curve] (3.6,0) circle (1.20);
  \foreach \a in {0,45,...,315}{\draw[component] ({3.6+0.66*cos(\a)},{0.66*sin(\a)})--({3.6+1.12*cos(\a)},{1.12*sin(\a)});}
  \node at (3.6,0) {$a$}; \node[above right] at (4.35,0.75) {$b$};
  \node[TextGray] at (3.6,-1.45) {同轴圆柱截面};
\end{tikzpicture}

\figcap{平行板、同心球与同轴圆柱电容器}\label{appE:fig:s21-legacy-02}
\end{center}}

\newcommand{\seriesParallelDiagram}{%
\begin{center}
\begin{tikzpicture}[scale=0.86,>=Latex]
  % parallel
  \draw[thick] (-5.0,1.25)--(-4.2,1.25) (-5.0,-1.25)--(-4.2,-1.25);
  \draw[thick] (-4.2,1.25)--(-4.2,-1.25) (-1.5,1.25)--(-1.5,-1.25);
  \draw[thick] (-4.2,0.72)--(-3.3,0.72) (-2.4,0.72)--(-1.5,0.72);
  \draw[very thick,DeepBlue] (-3.3,0.38)--(-3.3,1.06) (-2.4,0.38)--(-2.4,1.06);
  \draw[thick] (-4.2,-0.72)--(-3.3,-0.72) (-2.4,-0.72)--(-1.5,-0.72);
  \draw[very thick,DeepBlue] (-3.3,-1.06)--(-3.3,-0.38) (-2.4,-1.06)--(-2.4,-0.38);
  \node at (-2.85,1.28) {$C_1$}; \node at (-2.85,-1.28) {$C_2$};
  \node[TextGray] at (-3.25,-1.85) {并联：电压相同};
  % series
  \draw[thick] (0.8,0)--(1.8,0) (2.55,0)--(3.45,0) (4.2,0)--(5.2,0);
  \draw[very thick,DeepBlue] (1.8,-0.55)--(1.8,0.55) (2.55,-0.55)--(2.55,0.55);
  \draw[very thick,DeepBlue] (3.45,-0.55)--(3.45,0.55) (4.2,-0.55)--(4.2,0.55);
  \node at (2.18,0.95) {$C_1$}; \node at (3.82,0.95) {$C_2$};
  \node[TextGray] at (3.0,-1.85) {串联：中间孤立导体净电荷不变};
\end{tikzpicture}

\figcap{电容器的并联与串联}\label{appE:fig:s21-legacy-03}
\end{center}}

\newcommand{\polarizationDiagram}{%
\begin{center}
\begin{tikzpicture}[scale=0.87,>=Latex]
  \fill[PaleBlue] (-3.3,-1.5) rectangle (3.3,1.5);
  \draw[TextGray,thick] (-3.3,-1.5) rectangle (3.3,1.5);
  \foreach \x in {-2.6,-1.3,0,1.3,2.6}{
    \foreach \y in {-0.85,0,0.85}{
      \node[Ink] at (\x-0.16,\y) {$-$}; \node[DeepBlue] at (\x+0.16,\y) {$+$};
    }
  }
  \foreach \y in {-1.05,-0.35,0.35,1.05}{
    \node[Ink] at (-3.55,\y) {$-$}; \node[DeepBlue] at (3.55,\y) {$+$};
  }
  \draw[vector] (-1.0,2.05)--(1.0,2.05) node[midway,above] {$\vect E$};
  \draw[vector] (-1.0,-2.05)--(1.0,-2.05) node[midway,below] {$\vect P$};
  \node[TextGray,align=left,anchor=west] at (4.25,0.45)
    {分子偶极矩沿外场排列\\两端出现束缚面电荷};
\end{tikzpicture}

\figcap{外场作用下电介质的极化与束缚面电荷}\label{appE:fig:s21-legacy-04}
\end{center}}

\newcommand{\reciprocitySphereDiagram}{%
\begin{center}
\begin{tikzpicture}[x=1cm,y=1cm,>=Latex]
  \draw[curve] (-2.5,0) circle (1.18);
  \fill (-2.5,0) circle (1.8pt) node[below=4pt] {$O$};
  \fill[DeepBlue] (2.5,0) circle (2.6pt) node[above=5pt] {$P$};
  \draw[<->] (-2.5,0.25)--node[above] {$r$}(2.5,0.25);
  \node[align=center] at (-2.5,-1.65) {状态 I：球带单位电荷\\状态 II：$P$ 点放电荷 $q$};
  \node[align=left,anchor=west,TextGray] at (4.0,0.45)
    {$\displaystyle \sum_i Q_i\varphi_i'=\sum_iQ_i'\varphi_i$\\[4pt]
     交换“源”与“测量位置”，\\倒易关系保持不变};
\end{tikzpicture}

\figcap{格林倒易定理：导体球与球外一点的两种静电状态}\label{appE:fig:s21-legacy-05}
\end{center}}

\newcommand{\reciprocityConductorsDiagram}{%
\begin{center}
\begin{tikzpicture}[x=1cm,y=1cm,>=Latex]
  \fill[PaleBlue] plot[smooth cycle,tension=0.8] coordinates
    {(-4.1,0.2) (-3.5,1.25) (-2.15,1.05) (-1.65,0) (-2.3,-1.1) (-3.7,-0.95)};
  \draw[DeepBlue,thick] plot[smooth cycle,tension=0.8] coordinates
    {(-4.1,0.2) (-3.5,1.25) (-2.15,1.05) (-1.65,0) (-2.3,-1.1) (-3.7,-0.95)};
  \fill[SoftGray] plot[smooth cycle,tension=0.8] coordinates
    {(4.05,0.15) (3.35,1.2) (2.0,0.9) (1.55,-0.15) (2.35,-1.2) (3.75,-0.85)};
  \draw[Ink,thick] plot[smooth cycle,tension=0.8] coordinates
    {(4.05,0.15) (3.35,1.2) (2.0,0.9) (1.55,-0.15) (2.35,-1.2) (3.75,-0.85)};
  \node[align=center] at (-2.85,0) {$Q_1$\\$\varphi_1$};
  \node[align=center] at (2.85,0) {$Q_2$\\$\varphi_2$};
  \draw[<->,TextGray] (-1.15,0)--node[above] {静电感应}(1.15,0);
  \node[align=center,TextGray] at (0,-1.7)
    {第二状态用 $Q_1',Q_2',\varphi_1',\varphi_2'$ 表示；导体几何保持不变};
\end{tikzpicture}

\figcap{格林倒易定理：两个任意导体的两种静电平衡状态}\label{appE:fig:s21-legacy-06}
\end{center}}

\newcommand{\dielectricExperimentDiagram}{%
\begin{center}
\begin{tikzpicture}[x=0.86cm,y=0.86cm,>=Latex]
  \foreach \shift/\lab/\needle in {-4.3/{未插入介质：$C_0,U_0$}/32,4.3/{插入介质：$C=\varepsilon_rC_0$}/-8}{
    \begin{scope}[shift={(\shift,0)}]
      \draw[Ink,very thick] (-1.15,1.0)--(-1.15,3.1);
      \draw[Ink,very thick] (1.15,1.0)--(1.15,3.1);
      \node[DeepBlue] at (-1.45,2.1) {$+$}; \node[Ink] at (1.45,2.1) {$-$};
      \ifdim\shift pt>0pt
        \fill[DeepBlue!18] (-0.95,1.15) rectangle (0.95,2.95);
        \draw[DeepBlue,thick] (-0.95,1.15) rectangle (0.95,2.95);
        \node[DeepBlue,rotate=90] at (0,2.05) {电介质};
      \fi
      \draw[Ink,thick] (-1.15,1.0)--(-1.15,0.45)--(-0.75,0.45);
      \draw[Ink,thick] (1.15,1.0)--(1.15,0.45)--(0.75,0.45);
      \draw[rounded corners=3pt,fill=SoftGray,draw=Ink,thick] (-0.75,-1.5) rectangle (0.75,0.55);
      \draw[Ink] (0,-0.25) circle (0.48);
      \draw[DeepBlue,very thick] (0,-0.25)--({0.38*cos(\needle)},{-0.25+0.38*sin(\needle)});
      \foreach \a in {20,50,...,160}{\draw[TextGray] ({0.38*cos(\a)},{-0.25+0.38*sin(\a)})--({0.45*cos(\a)},{-0.25+0.45*sin(\a)});}
      \node[align=center] at (0,-2.05) {\lab};
    \end{scope}
  }
  \draw[vector] (-1.15,0.2)--(1.15,0.2) node[midway,above] {插入介质};
\end{tikzpicture}

\figcap{保持电荷不变时，插入电介质使电势差减小、电容增大}\label{appE:fig:s21-legacy-07}
\end{center}}

\newcommand{\touchingSpheresDiagram}{%
\begin{center}
\begin{tikzpicture}[scale=0.92,>=Latex]
  \draw[curve] (-1.3,0) circle (1.3); \draw[curve] (1.3,0) circle (1.3);
  \fill (0,0) circle (1.4pt) node[below] {接触点};
  \draw[<->] (1.3,0)--node[above] {$R$}(2.6,0);
  \foreach \x/\r in {-1.75/2.3,-0.95/1.8,-0.48/1.25,0.48/1.25,0.95/1.8,1.75/2.3}{
    \fill[DeepBlue] (\x,0) circle (\r pt);
  }
  \node[TextGray] at (0,-2.05) {接触极限可由双球坐标或无穷像电荷序列求得};
\end{tikzpicture}

\figcap{两个等半径相切球连成一个等势导体}\label{appE:fig:s21-legacy-08}
\end{center}}

\newcommand{\twoCylinderDiagram}{%
\begin{center}
\begin{tikzpicture}[scale=0.9,>=Latex]
  \draw[curve] (-2.2,0) circle (0.85); \draw[curve] (2.2,0) circle (0.85);
  \node at (-2.2,0) {$+\lambda$}; \node at (2.2,0) {$-\lambda$};
  \fill ( -2.2,0) circle (1.3pt); \fill (2.2,0) circle (1.3pt);
  \draw[<->] (-2.2,-1.28)--node[fill=white,inner sep=1pt] {$L$}(2.2,-1.28);
  \draw[<->] (-2.2,0)--node[above] {$R$}(-1.35,0);
  \draw[DeepBlue,thick] (-1.55,0.8) .. controls (0,2.0) .. (1.55,0.8);
  \draw[DeepBlue,thick] (-1.55,-0.8) .. controls (0,-2.0) .. (1.55,-0.8);
  \draw[DeepBlue,thick] (-1.35,0.25) .. controls (0,0.78) .. (1.35,0.25);
  \draw[DeepBlue,thick] (-1.35,-0.25) .. controls (0,-0.78) .. (1.35,-0.25);
  \node[TextGray] at (0,-2.35) {$\cosh\tau_0=L/(2R)$，两圆柱面对应 $\tau=\pm\tau_0$};
\end{tikzpicture}

\figcap{两根等半径平行圆柱导体的横截面}\label{appE:fig:s21-legacy-09}
\end{center}}

\newcommand{\splitDielectricDiagram}{%
\begin{center}
\begin{tikzpicture}[scale=0.83,>=Latex]
  % side-by-side
  \draw[DeepBlue,very thick] (-5.1,1.0)--(-1.9,1.0) (-5.1,-1.0)--(-1.9,-1.0);
  \fill[PaleBlue] (-5.1,-0.95) rectangle (-3.5,0.95);
  \draw[guide] (-3.5,-1.0)--(-3.5,1.0);
  \node at (-4.3,0) {$\varepsilon$}; \node at (-2.7,0) {$\varepsilon_0$};
  \node[TextGray] at (-3.5,-1.55) {按面积分割：并联};
  % layered
  \draw[DeepBlue,very thick] (1.2,1.0)--(4.4,1.0) (1.2,-1.0)--(4.4,-1.0);
  \fill[PaleBlue] (1.2,0) rectangle (4.4,0.95);
  \draw[guide] (1.2,0)--(4.4,0);
  \node at (2.8,0.5) {$\varepsilon$}; \node at (2.8,-0.5) {$\varepsilon_0$};
  \node[TextGray] at (2.8,-1.55) {按厚度分割：串联};
\end{tikzpicture}

\figcap{部分填充介质的两种等效方式}\label{appE:fig:s21-legacy-10}
\end{center}}

\newcommand{\dielectricPullDiagram}{%
\begin{center}
\begin{tikzpicture}[scale=0.88,>=Latex]
  \draw[DeepBlue,very thick] (-3.8,0.8)--(3.8,0.8) (-3.8,-0.8)--(3.8,-0.8);
  \fill[PaleBlue] (-3.8,-0.72) rectangle (0.6,0.72);
  \draw[TextGray,thick] (-3.8,-0.72) rectangle (0.6,0.72);
  \draw[<->] (-3.8,-1.22)--node[fill=white,inner sep=1pt] {$x$}(0.6,-1.22);
  \draw[<->] (0.6,-1.22)--node[fill=white,inner sep=1pt] {$a-x$}(3.8,-1.22);
  \draw[<->] (4.15,-0.8)--node[right] {$d$}(4.15,0.8);
  \draw[vector] (0.55,0)--(2.2,0) node[midway,above] {$\vect F$};
  \node[TextGray] at (0,-1.95) {介质进入越深，等效电容越大，电场把介质继续拉入};
\end{tikzpicture}

\figcap{部分插入平行板电容器的介质板}\label{appE:fig:s21-legacy-11}
\end{center}}

\newcommand{\dielectricSphereDiagram}{%
\begin{center}
\begin{tikzpicture}[scale=0.88,>=Latex]
  \fill[PaleBlue] (0,0) circle (1.35); \draw[curve] (0,0) circle (1.35);
  \foreach \a in {-60,-30,0,30,60}{
    \node[Ink] at ({-1.08*cos(\a)},{1.08*sin(\a)}) {$-$};
    \node[DeepBlue] at ({1.08*cos(\a)},{1.08*sin(\a)}) {$+$};
  }
  \draw[vector] (-4.3,1.95)--(-2.8,1.95) node[midway,above] {$\vect E_0$};
  \draw[DeepBlue,thick,-{Latex[length=2mm]}] (-4.2,1.72)--(4.2,1.72);
  \draw[DeepBlue,thick,-{Latex[length=2mm]}] (-4.2,-1.72)--(4.2,-1.72);
  \draw[DeepBlue,thick] (-4.2,0.92) .. controls (-2.6,0.92) and (-1.8,0.62) .. (-1.26,0.48);
  \draw[DeepBlue,thick,-{Latex[length=2mm]}] (-1.26,0.48)--(1.26,0.48);
  \draw[DeepBlue,thick,-{Latex[length=2mm]}] (1.26,0.48) .. controls (1.8,0.62) and (2.6,0.92) .. (4.2,0.92);
  \draw[DeepBlue,thick] (-4.2,-0.92) .. controls (-2.6,-0.92) and (-1.8,-0.62) .. (-1.26,-0.48);
  \draw[DeepBlue,thick,-{Latex[length=2mm]}] (-1.26,-0.48)--(1.26,-0.48);
  \draw[DeepBlue,thick,-{Latex[length=2mm]}] (1.26,-0.48) .. controls (1.8,-0.62) and (2.6,-0.92) .. (4.2,-0.92);
  \draw[DeepBlue,thick] (-4.2,0)--(-1.35,0);
  \draw[DeepBlue,thick,-{Latex[length=2mm]}] (-1.35,0)--(1.35,0);
  \draw[DeepBlue,thick,-{Latex[length=2mm]}] (1.35,0)--(4.2,0);
  \node at (0,-0.82) {$\varepsilon$};
  \node[TextGray] at (0,-2.35) {$\vect E_{\rm in}=\dfrac{3\varepsilon_0}{\varepsilon+2\varepsilon_0}\vect E_0$};
\end{tikzpicture}

\figcap{匀强电场中的线性电介质球}\label{appE:fig:s21-legacy-12}
\end{center}}

\newcommand{\chargingDiagram}{%
\begin{center}
\begin{tikzpicture}[scale=0.9,>=Latex]
  \draw[thick] (-3.8,0)--(-2.9,0);
  \draw[thick] (-2.9,-0.55)--(-2.9,0.55); \draw[very thick,DeepBlue] (-2.55,-0.35)--(-2.55,0.35);
  \node[below] at (-2.72,-0.62) {$\mathcal E$};
  \draw[thick] (-2.55,0)--(-1.7,0);
  \draw[thick,decorate,decoration={zigzag,segment length=6pt,amplitude=2.5pt}] (-1.7,0)--(0.2,0);
  \node[above] at (-0.75,0.35) {$R$};
  \draw[thick] (0.2,0)--(1.25,0);
  \draw[very thick,DeepBlue] (1.25,-0.55)--(1.25,0.55) (1.85,-0.55)--(1.85,0.55);
  \node[above] at (1.55,0.75) {$C$};
  \draw[thick] (1.85,0)--(2.8,0)--(2.8,-1.35)--(-3.8,-1.35)--(-3.8,0);
  \node[TextGray,align=left,anchor=west] at (4.0,-0.1)
    {电源功 $C\mathcal E^2$\\电场能 $C\mathcal E^2/2$\\焦耳热 $C\mathcal E^2/2$};
\end{tikzpicture}

\figcap{电阻参与的电容器充电过程}\label{appE:fig:s21-legacy-13}
\end{center}}

\newcommand{\starCapacitorDiagram}{%
\begin{center}
\begin{tikzpicture}[scale=0.94,>=Latex]
  \fill[DeepBlue] (0,0) circle (2pt) node[below right] {$\varphi_0$};
  \foreach \a/\c/\p in {90/C_1/\varphi_1,210/C_2/\varphi_2,330/C_3/\varphi_3}{
    \draw[thick] (0,0)--({0.95*cos(\a)},{0.95*sin(\a)});
    \draw[very thick,DeepBlue] ({0.78*cos(\a)-0.35*sin(\a)},{0.78*sin(\a)+0.35*cos(\a)})--({0.78*cos(\a)+0.35*sin(\a)},{0.78*sin(\a)-0.35*cos(\a)});
    \draw[very thick,DeepBlue] ({1.12*cos(\a)-0.35*sin(\a)},{1.12*sin(\a)+0.35*cos(\a)})--({1.12*cos(\a)+0.35*sin(\a)},{1.12*sin(\a)-0.35*cos(\a)});
    \draw[thick] ({1.12*cos(\a)},{1.12*sin(\a)})--({2.2*cos(\a)},{2.2*sin(\a)});
    \node at ({1.48*cos(\a)+0.45*sin(\a)},{1.48*sin(\a)-0.45*cos(\a)}) {$\c$};
    \node at ({2.55*cos(\a)},{2.55*sin(\a)}) {$\p$};
  }
  \node[TextGray,align=left,anchor=west] at (4.0,0.45)
    {中心导体孤立且初始中性\\$\sum C_i(\varphi_0-\varphi_i)=0$};
\end{tikzpicture}

\figcap{三电容星形连接的浮置中心节点}\label{appE:fig:s21-legacy-14}
\end{center}}

\newcommand{\atmosphereCapacitorDiagram}{%
\begin{center}
\begin{tikzpicture}[scale=0.83,>=Latex]
  \fill[PaleBlue] (0,0) circle (2.18); \fill[white] (0,0) circle (1.86);
  \fill[SoftGray] (0,0) circle (1.62);
  \draw[curve] (0,0) circle (1.62); \draw[DeepBlue,thick] (0,0) circle (2.18);
  \node at (0,0) {地球};
  \foreach \a in {0,35,...,325}{
    \node[Ink] at ({1.72*cos(\a)},{1.72*sin(\a)}) {$-$};
    \node[DeepBlue] at ({2.05*cos(\a)},{2.05*sin(\a)}) {$+$};
  }
  \draw[<->] (2.35,0)--node[right] {$h$}(2.35,0.56);
  \node[TextGray,align=left,anchor=west] at (4.1,0.7)
    {内半径 $R_\oplus$\\外半径 $R_\oplus+h$\\大气同时提供电容与漏电电阻};
\end{tikzpicture}

\figcap{理想化的大气球形电容器模型}\label{appE:fig:s21-legacy-15}
\end{center}}


\begin{origexamplebox}{【统统算出来】}\label{appE:ex:s21-32-1}
求下列三种电容器的电容：
\begin{enumerate}[label=(\arabic*),leftmargin=2.2em]
  \item 面积 $S=l^2$、间距 $d$ 的平行板电容器，$l\gg d$；
  \item 内、外半径为 $R_1,R_2$ 的同心球形电容器，$R_1<R_2$；
  \item 内、外半径为 $a,b$、长度为 $l$ 的同轴圆柱形电容器，$l\gg b$。
\end{enumerate}

\QAENamed{\QAECapacitors{0}}{平行板、同心球与同轴圆柱三种电容器的尺寸；后两图均为中心横截面。}{quality:appE:q6}
\end{origexamplebox}

\begin{xesolution}
三题都让内、外电极分别带 $+Q$、$-Q$，先由高斯定律求两极间电场，再积分求电压。

\textbf{(1) 平行板。}忽略边缘效应，板间场为
\[
  E=\frac{Q}{\varepsilon_0S},\qquad
  U=Ed=\frac{Qd}{\varepsilon_0S}.
\]
因此
\begin{equation}
  \boxed{C_{\rm pp}=\frac QU=\frac{\varepsilon_0S}{d}}.
\end{equation}

\textbf{(2) 同心球。}在 $R_1<r<R_2$ 内，
\[
  E(r)4\pi r^2=\frac Q{\varepsilon_0},
  \qquad E(r)=\frac{Q}{4\pi\varepsilon_0r^2}.
\]
故
\[
  U=\int_{R_1}^{R_2}E(r)\dd r
  =\frac{Q}{4\pi\varepsilon_0}\left(\frac1{R_1}-\frac1{R_2}\right),
\]
从而
\begin{equation}
  \boxed{C_{\rm sph}=4\pi\varepsilon_0
  \frac{R_1R_2}{R_2-R_1}}.
\end{equation}
令 $R_2\to\infty$，得到孤立球的 $C=4\pi\varepsilon_0R_1$。

\textbf{(3) 同轴圆柱。}柱面高斯面给出
\[
  E(r)2\pi rl=\frac Q{\varepsilon_0},
  \qquad E(r)=\frac{Q}{2\pi\varepsilon_0lr}.
\]
所以
\[
  U=\int_a^bE(r)\dd r=\frac{Q}{2\pi\varepsilon_0l}\ln\frac ba,
\]
以及
\begin{equation}
  \boxed{C_{\rm cyl}=\frac{2\pi\varepsilon_0l}{\ln(b/a)}}.
\end{equation}
三种推导中的近似只用于忽略开口或极板边缘的非均匀场；在理想同心球中不存在边缘近似。

\QAENamed{\QAECapacitors{1}}{积分路径沿两极间的场线；同心球与同轴柱使用不同的高斯面积。}{quality:appE:a6}

图~\ref{quality:appE:a6}内外电极之间分别使用平面、球面和柱面对称性。球面面积随 $r^2$ 增长，柱面面积随 $r$ 增长，因而径向积分一个给倒数差、一个给对数；不能仅凭外形把两式套换。
\end{xesolution}

\begin{origexamplebox}{【改造电容器】}\label{appE:ex:s21-32-4}
平行板电容器的极板面积为 $S$、间距为 $d$。
\begin{enumerate}[label=(\arabic*),leftmargin=2.2em]
  \item 沿面积方向，在左半边完全填入相对介电常数为 $\varepsilon_r>1$ 的介质；
  \item 沿厚度方向，在上半层填入相对介电常数为 $\varepsilon_r$ 的介质。
\end{enumerate}
分别求改造后的电容，并判断增减。

忽略边缘效应对总电容的修正，介质为均匀线性电介质。

\QAENamed{\QAESplitQuestion}{半面积填充与半厚度填充；蓝色表示介质，白色表示真空。}{quality:appE:q7}
\end{origexamplebox}

\begin{xesolution}
\textbf{(1) 按面积分成左右两区。}两区具有相同板间距与电压，相当于并联：
\[
  C_1=\frac{\varepsilon_r\varepsilon_0(S/2)}d,
  \qquad C_2=\frac{\varepsilon_0(S/2)}d.
\]
故
\begin{equation}
  \boxed{C_{\rm left}=C_1+C_2
  =\frac{\varepsilon_0S}{2d}(\varepsilon_r+1)}.
\end{equation}

\textbf{(2) 按厚度分成上下两层。}两层面积相同，自由面电荷密度相同，电位移的法向分量连续，
相当于串联：
\[
  C_1=\frac{\varepsilon_r\varepsilon_0S}{d/2},
  \qquad C_2=\frac{\varepsilon_0S}{d/2}.
\]
于是
\begin{equation}
  \boxed{C_{\rm upper}=\left(\frac1{C_1}+\frac1{C_2}\right)^{-1}
  =\frac{2\varepsilon_r}{\varepsilon_r+1}\frac{\varepsilon_0S}{d}}.
\end{equation}
因 $\varepsilon_r>1$，两种改造都使电容变大；但一般有
\[
  C_{\rm left}-C_{\rm upper}
  =\frac{\varepsilon_0S}{d}\frac{(\varepsilon_r-1)^2}{2(\varepsilon_r+1)}\ge0,
\]
所以相同介质量按面积并排放置，增容更明显。

\QAENamed{\splitDielectricDiagram}{两区域等电压形成并联，两层等法向电位移形成串联；$\varepsilon=\varepsilon_r\varepsilon_0$。}{quality:appE:a7}

图~\ref{quality:appE:a7}左图沿面积切开，两块介质都连接同一对电极；右图沿厚度切开，中间界面没有自由面电荷，故两层法向 $D$ 相同。这些边界条件分别决定并联和串联，而不是按介质块数判断。
\end{xesolution}

\begin{origexamplebox}{【奇妙的吸引力】}\label{appE:ex:s21-32-5}
小明又需要你的帮助了。他把一块相对介电常数为 $\varepsilon_r>1$ 的电介质插入平行板电容器中。
极板面积为 $S=ab$、间距为 $d$，介质厚度等于板间距，并沿边长 $a$ 的方向插入深度 $x$。
电介质会受到继续向电容器内部运动的吸引力；这个力从微观上怎样产生并不容易直接看清，请利用虚功法
把它算出来。
\begin{enumerate}[label=(\arabic*),leftmargin=2.2em]
  \item 电容器电荷 $Q$ 保持不变，用介质内电场 $E$ 表示吸入力；
  \item 电容器接在恒压电源上，用介质内电场 $E$ 表示吸入力。
\end{enumerate}

忽略边缘效应对总电容的修正，介质为均匀线性电介质。

\QAENamed{\QAEInsertQuestion}{介质沿 $+x$ 方向进入，插入长度 $0<x<a$；图中尚未给出受力。}{quality:appE:q8}
\end{origexamplebox}

\begin{xesolution}
在薄间隙、介质前缘远离极板两端的内部插入区间，忽略边缘效应对总电容的高阶修正，已填充区域与真空区域相当于并联：
\begin{equation}
  C(x)=\frac{\varepsilon_0b}{d}\,[a+(\varepsilon_r-1)x],
  \qquad
  \frac{\dd C}{\dd x}=\frac{\varepsilon_0b}{d}(\varepsilon_r-1).
\end{equation}

\textbf{(1) 固定电荷。}系统静电能 $W=Q^2/[2C(x)]$，介质所受力沿 $x$ 增大方向：
\[
  F=-\left(\frac{\partial W}{\partial x}\right)_Q
  =\frac{Q^2}{2C^2}\frac{\dd C}{\dd x}
  =\frac{U^2}{2}\frac{\dd C}{\dd x}.
\]
介质内场 $E=U/d$，故
\begin{equation}
  \boxed{F_Q=\frac12\varepsilon_0(\varepsilon_r-1)bdE^2}.
\end{equation}

\textbf{(2) 固定电压。}电容器能量 $CU^2/2$ 随 $C$ 增大而增大，但电源还会输送电荷并做功
$U\dd Q=U^2\dd C$。把电源计入后，机械势能的增量是
$\dd(CU^2/2)-U\dd Q=-U^2\dd C/2$，所以
\[
  F=\frac{U^2}{2}\frac{\dd C}{\dd x}.
\]
仍用 $E=U/d$ 表示，得到
\begin{equation}
  \boxed{F_U=\frac12\varepsilon_0(\varepsilon_r-1)bdE^2}.
\end{equation}
两式形式相同，是因为题目要求用各自状态下的 $E$ 表示；固定 $Q$ 时 $E$ 随 $x$ 改变，固定
$U$ 时 $E$ 不变。力都把介质拉入电容器，因为介质进入会增大电容。

\QAENamed{\dielectricPullDiagram}{吸引力指向插入深度增加方向；固定电压时须连同电源功计算。}{quality:appE:a8}

图~\ref{quality:appE:a8}的箭头沿介质右端向电容器内部，即 $x$ 增大方向。力来自插入前缘附近的非均匀边缘场；忽略边缘场对总电容的高阶修正，并不意味着可以在微观受力分析中删去边缘场。虚功法用 $C'(x)$ 保留了该作用的整体效果。
\end{xesolution}

\begin{origexamplebox}{【充电过程】}\label{appE:ex:s21-32-8}
用电动势为 $\mathcal E$ 的理想电源给初始未充电的电容 $C$ 充电。
\begin{enumerate}[label=(\arabic*),leftmargin=2.2em]
  \item 电路中串联电阻 $R$，充电完成前电源做功和电阻发热各为多少；
  \item 令 $R$ 趋于零，两个结果怎样变化；
  \item 若理想电路中完全没有电阻，应如何理解？
\end{enumerate}

\QAENamed{\QAERCQuestion}{初始未充电的电容器与串联电阻、电源；讨论闭合后的过程。}{quality:appE:q9}
\end{origexamplebox}

\begin{xesolution}
\textbf{(1)} 充电完成后 $Q=C\mathcal E$。电源搬运总电荷 $Q$，故
\begin{equation}
  \boxed{W_{\rm source}=\mathcal EQ=C\mathcal E^2}.
\end{equation}
电容器最终储能
\[
  W_C=\frac12C\mathcal E^2.
\]
其余能量变成焦耳热：
\begin{equation}
  \boxed{Q_R=W_{\rm source}-W_C=\frac12C\mathcal E^2}.
\end{equation}
也可由 $i(t)=\mathcal E e^{-t/(RC)}/R$ 积分验证：
$\int_0^\infty i^2R\dd t=C\mathcal E^2/2$。

\textbf{(2)} 当 $R\to0^+$，充电时间常数 $RC\to0$、初始电流峰值趋于无穷，但热量积分仍为
$C\mathcal E^2/2$，与 $R$ 的数值无关。这是先对有限 $R$ 完成积分、再取极限的结果。

\textbf{(3)} 完全没有电阻、电感和辐射的理想回路不能由一阶 $RC$ 方程描述；它要求电容电压在
零时间内突变，产生无穷大电流，数学模型自相矛盾。真实电路总有寄生电感、导线电阻和电磁辐射。
若保留电感而不含耗散，能量会在电容电场与电感磁场之间振荡，电容电压并不单调趋于
$\mathcal E$；微小电阻或辐射最终耗散那一半能量。因而“零电阻”只能作为 $R\to0^+$ 的极限，
不能把所有寄生参数同时删掉后仍期待同样的稳态充电过程。

\QAENamed{\chargingDiagram}{有限电阻充电的三项能量账；先积分，再取电阻趋零极限。}{quality:appE:a9}

图~\ref{quality:appE:a9}将电源功拆成电场储能与焦耳热。初始未充电是两项各占一半的前提；若电容预先有电压，必须重新计算搬运的电荷和初末储能，不能仍套用这个数值。
\end{xesolution}

\begin{origexamplebox}{【超大电容】}\label{appE:ex:s21-32-11}
考虑地球与大气的理想化电容器模型：地球表面是其中一个极板，带负电荷；
高空的大气等效为另一个极板，带正电荷。把地球表面和高空 $h=\SI{5}{km}$ 处的大气层近似为
同心球形电容器。取地球半径
$R_\oplus=\SI{6.4e6}{m}$，地球带电量绝对值 $Q=\SI{5e5}{C}$，空气电阻率
$\rho=\SI{3e13}{\ohm.m}$。
\begin{enumerate}[label=(\arabic*),leftmargin=2.2em]
  \item 求球形电容器的电容；
  \item 求地表电场能量密度和电容器总能量；
  \item 求大气层的径向漏电电阻；
  \item 求 $RC$ 放电时间常数；
  \item 若每次雷电平均向地面补充 $\SI{25}{C}$ 电荷，维持该电荷量每天约需多少次雷电？
\end{enumerate}

\QAENamed{\QAEEarthQuestion}{地表与等效大气外壳的同心球模型；图中的大气层厚度已放大。}{quality:appE:q10}
\end{origexamplebox}

\begin{xesolution}
令内、外半径分别为 $R=R_\oplus$、$R+h$。

\textbf{(1) 电容。}
\begin{equation}
  C=4\pi\varepsilon_0\frac{R(R+h)}h
  \approx\boxed{\SI{0.912}{F}}.
\end{equation}
因 $h\ll R$，也可用薄层近似 $C\approx4\pi\varepsilon_0R^2/h$。

\textbf{(2) 地表场与能量。}地表外侧
\[
  E(R)=\frac{Q}{4\pi\varepsilon_0R^2}\approx\SI{1.10e2}{V.m^{-1}}.
\]
故
\begin{equation}
  \boxed{u(R)=\frac12\varepsilon_0E^2\approx\SI{5.33e-8}{J.m^{-3}}}.
\end{equation}
总能量
\begin{equation}
  \boxed{W=\frac{Q^2}{2C}\approx\SI{1.37e11}{J}}.
\end{equation}

\textbf{(3) 径向电阻。}半径 $r$ 处厚度 $\dd r$ 的球壳电阻
$\dd R_{\rm el}=\rho\dd r/(4\pi r^2)$，所以
\begin{equation}
  R_{\rm el}=\frac\rho{4\pi}\left(\frac1R-\frac1{R+h}\right)
  \approx\boxed{\SI{2.91e2}{\ohm}}.
\end{equation}

\textbf{(4) 时间常数。}相乘时几何因子恰好消去：
\begin{equation}
  \boxed{\tau=R_{\rm el}C=\rho\varepsilon_0
  \approx\SI{2.66e2}{s}\approx\SI{4.43}{min}}.
\end{equation}

\textbf{(5) 雷电补充率。}漏电电流的平均量级为
\[
  I=\frac Q\tau\approx\SI{1.88e3}{A}.
\]
每天漏失电荷 $I\times\SI{86400}{s}$，所需雷电次数
\begin{equation}
  \boxed{N=\frac{Q}{\tau}\frac{\SI{86400}{s}}{\SI{25}{C}}
  \approx6.5\times10^6\ \text{次/天}}.
\end{equation}
这是题设参数下的数量级模型；真实全球电路的电荷分布、天气电导率和雷暴补充机制更复杂。

\QAENamed{\QAEEarthQuestion}{同一球壳面积 $4\pi r^2$ 分别进入场、电容和漏电电阻，形成几何因子的抵消。}{quality:appE:a10}

图~\ref{quality:appE:a10}中的径向薄壳具有面积 $4\pi r^2$。电阻积分类似串联薄层，电势积分也沿相同半径进行，故 $R_{\rm el}C=\rho\varepsilon_0$ 与球壳的两端半径无关；这要求电阻率在整个壳层内均匀。
\end{xesolution}

\endgroup

\Needspace{18\baselineskip}

\section{恒定电流}

\begingroup

\graphicspath{{../../../sample/xes-original/twenty-second-section/}{../../../sample/xes-original/twenty-second-section/assets/}}

% Restrained TikZ redraws for Lesson 33.

\newcommand{\currentDensityDiagram}{%
\begin{center}
\begin{tikzpicture}[scale=0.88,>=Latex]
  \fill[PaleBlue] (-3.6,-1.15) rectangle (3.6,1.15);
  \draw[curve] (-3.6,-1.15) rectangle (3.6,1.15);
  \draw[DeepBlue,very thick] (0,-1.15)--(0,1.15);
  \foreach \y in {-0.75,-0.25,0.25,0.75}{
    \draw[vector] (-2.8,\y)--(-0.3,\y);
    \draw[vector] (0.3,\y)--(2.8,\y);
  }
  \node[above] at (0,1.2) {截面 $S$};
  \node[TextGray] at (0,-1.7) {$I=\displaystyle\int_S\vect j\cdot\dd\vect S$，箭头表示约定的电流方向};
\end{tikzpicture}

\figcap{补图 33-A\quad 电流是电流密度穿过截面的通量}\label{appE:fig:s22-legacy-01}
\end{center}}

\newcommand{\ivCharacteristicDiagram}{%
\begin{center}
\begin{tikzpicture}[scale=0.88,>=Latex]
  \draw[axis] (-3.5,0)--(3.7,0) node[right] {$U$};
  \draw[axis] (0,-2.6)--(0,2.8) node[above] {$I$};
  \draw[DeepBlue,very thick] (-2.7,-2.1)--(2.7,2.1) node[above right] {理想电阻};
  \draw[Ink,very thick] (-2.9,-2.15) .. controls (-1.8,-0.45) and (-0.7,-0.25) .. (0,0)
    .. controls (0.7,0.25) and (1.8,0.45) .. (2.9,2.15);
  \node[Ink,anchor=west] at (2.35,1.2) {非欧姆元件};
  \draw[guide] (2.0,0)--(2.0,1.56);
  \node[TextGray] at (0,-3.15) {选定端口电压、电流正方向后，曲线完整描述二端元件的外特性};
\end{tikzpicture}

\figcap{二端元件及其伏安特性曲线}\label{appE:fig:s22-legacy-02}
\end{center}}

\newcommand{\ohmHistoryContextDiagram}{%
\begin{center}
\begin{tikzpicture}[scale=0.82,>=Latex]
  \node[inner sep=0] at (-3.7,0)
    {\includegraphics[height=3.4cm]{ohm-portrait-redraw.png}};
  \node[TextGray] at (-3.7,-2.05) {格奥尔格·西蒙·欧姆};
  \begin{scope}[xshift=2.0cm]
    \draw[thick,rounded corners=4pt] (-2.3,-1.45) rectangle (3.0,1.45);
    \foreach \y in {-1.2,-0.8,...,1.2}{
      \draw[TextGray] (-2.42,\y)--(-2.3,\y);
      \draw[TextGray] (3.0,\y)--(3.12,\y);
    }
    \draw[thick] (-1.6,0)--(-0.85,0);
    \draw[fill=SoftGray,thick,rounded corners=2pt] (-0.85,-0.38) rectangle (1.5,0.38);
    \draw[thick] (1.5,0)--(2.25,0);
    \foreach \x in {-0.45,0.05,0.55,1.05}{\draw[DeepBlue,very thick] (\x,-0.38)--(\x,0.38);}
    \node[DeepBlue,font=\large] at (0.3,0.95) {$U=IR$};
    \node[TextGray,font=\small] at (0.3,-0.95) {色环电阻与欧姆定律};
  \end{scope}
\end{tikzpicture}
\figcap{欧姆肖像与纪念欧姆定律的邮票图案}\label{appE:fig:s22-legacy-03}
\end{center}}

\newcommand{\uniformResistorDiagram}{%
\begin{center}
\begin{tikzpicture}[x=1cm,y=1cm,>=Latex]
  \fill[SoftGray] (-3.2,-0.72) rectangle (3.2,0.72);
  \draw[Ink,thick] (-3.2,-0.72) rectangle (3.2,0.72);
  \draw[Ink,thick] (-3.2,0) ellipse (0.24 and 0.72);
  \draw[Ink,thick] (3.2,0) ellipse (0.24 and 0.72);
  \draw[vector] (-2.5,1.25)--(2.5,1.25) node[midway,above] {$I$};
  \draw[<->] (-3.2,-1.2)--node[below] {$l$}(3.2,-1.2);
  \node at (-3.85,0.25) {$A$};
  \node[TextGray,align=left,anchor=west] at (4.3,0.35)
    {均匀材料：电阻率 $\rho$\\截面积 $A$，长度 $l$\\$\displaystyle R=\rho\frac lA$};
\end{tikzpicture}

\figcap{均匀圆柱形电阻的几何参数}\label{appE:fig:s22-legacy-04}
\end{center}}

\newcommand{\gasDischargeContextDiagram}{%
\begin{center}
\begin{tikzpicture}[x=1.03cm,y=0.92cm,>=Latex]
  \draw[axis] (0,0)--(8.1,0) node[right] {$I$};
  \draw[axis] (0,0)--(0,4.15) node[above] {$U$};
  \foreach \x/\lab in {0.8/$10^{-15}$,2.5/$10^{-10}$,4.5/$10^{-5}$,6.4/$10^0$}{
    \draw (\x,0.08)--(\x,-0.08) node[below] {\lab};
  }
  \foreach \y/\lab in {1.0/300,2.0/600,3.35/1000}{
    \draw (0.08,\y)--(-0.08,\y) node[left] {\lab};
  }
  \draw[guide,dashed] (4.25,0)--(4.25,3.8);
  \draw[guide,dashed] (6.25,0)--(6.25,3.8);
  \node[TextGray] at (2.1,3.75) {暗放电};
  \node[TextGray] at (5.2,3.75) {辉光放电};
  \node[TextGray] at (7.1,3.75) {弧光放电};
  \draw[DeepBlue,very thick]
    (0.45,0.25) .. controls (0.7,0.28) and (0.82,0.55) .. (0.92,1.75)
    .. controls (1.02,2.05) and (1.18,2.12) .. (1.35,2.15)
    -- (3.85,2.3)
    .. controls (4.35,2.28) and (4.5,2.05) .. (4.8,1.75)
    .. controls (5.1,1.35) and (5.4,0.9) .. (5.75,0.95)
    .. controls (6.05,1.0) and (6.1,1.75) .. (6.22,2.25)
    .. controls (6.3,2.75) and (6.36,3.45) .. (6.45,3.5)
    .. controls (6.55,3.3) and (6.42,0.6) .. (6.75,0.42)
    .. controls (7.1,0.28) and (7.55,0.48) .. (7.85,0.58);
  \fill[DeepBlue] (3.85,2.3) circle (2.4pt);
  \fill[DeepBlue] (5.75,0.95) circle (2.4pt);
  \fill[DeepBlue] (6.45,3.5) circle (2.4pt);
  \fill[DeepBlue] (6.75,0.42) circle (2.4pt);
  \draw[<-,Ink,thick] (5.25,1.05)--(4.55,0.45)
    node[below left,align=center] {$\dd U/\dd I<0$\\负微分电阻};
\end{tikzpicture}

\figcap{典型气体放电的电压—电流曲线；横轴按对数刻度标示}\label{appE:fig:s22-legacy-05}
\end{center}}

\newcommand{\drudeDiagram}{%
\begin{center}
\begin{tikzpicture}[scale=0.85,>=Latex]
  \fill[SoftGray] (-4.0,-1.35) rectangle (4.0,1.35);
  \foreach \x in {-3.4,-2.3,...,3.3}{\foreach \y in {-0.85,0,0.85}{\fill[TextGray] (\x,\y) circle (2.4pt);}}
  \draw[DeepBlue,thick,-{Latex[length=2mm]}] (-3.5,0.55) .. controls (-2.7,1.05) and (-2.1,-0.8) .. (-1.2,-0.25)
    .. controls (-0.4,0.45) and (0.5,-0.8) .. (1.25,0.1)
    .. controls (1.9,0.85) and (2.65,-0.45) .. (3.45,0.35);
  \draw[vector] (-1.1,2.0)--(1.1,2.0) node[midway,above] {$\vect E$};
  \draw[vector] (1.0,-1.95)--(-1.0,-1.95) node[midway,below] {电子平均漂移};
  \node[TextGray] at (0,-2.55) {高速无规则运动叠加很小的定向漂移；碰撞使动量不断弛豫};
\end{tikzpicture}

\figcap{导体内电子的高速无规则运动与微小定向漂移}\label{appE:fig:s22-legacy-06}
\end{center}}

\newcommand{\resistorConnectionDiagram}{%
\begin{center}
\begin{tikzpicture}[scale=0.85,>=Latex]
  \draw[thick] (-5.0,0)--(-4.2,0);
  \draw[thick,decorate,decoration={zigzag,segment length=6pt,amplitude=2.5pt}] (-4.2,0)--(-2.8,0);
  \draw[thick] (-2.8,0)--(-2.2,0);
  \draw[thick,decorate,decoration={zigzag,segment length=6pt,amplitude=2.5pt}] (-2.2,0)--(-0.8,0);
  \draw[thick] (-0.8,0)--(-0.1,0);
  \node[above] at (-3.5,0.45) {$R_1$}; \node[above] at (-1.5,0.45) {$R_2$};
  \node[TextGray] at (-2.55,-1.0) {串联：电流相同};
  \draw[thick] (1.1,1.05)--(1.8,1.05) (1.1,-1.05)--(1.8,-1.05);
  \draw[thick] (1.8,1.05)--(1.8,-1.05) (5.2,1.05)--(5.2,-1.05);
  \draw[thick] (1.8,0.62)--(2.45,0.62);
  \draw[thick,decorate,decoration={zigzag,segment length=6pt,amplitude=2.5pt}] (2.45,0.62)--(4.5,0.62);
  \draw[thick] (4.5,0.62)--(5.2,0.62);
  \draw[thick] (1.8,-0.62)--(2.45,-0.62);
  \draw[thick,decorate,decoration={zigzag,segment length=6pt,amplitude=2.5pt}] (2.45,-0.62)--(4.5,-0.62);
  \draw[thick] (4.5,-0.62)--(5.2,-0.62);
  \node[above] at (3.48,0.88) {$R_1$}; \node[below] at (3.48,-0.88) {$R_2$};
  \node[TextGray] at (3.15,-1.65) {并联：电压相同};
\end{tikzpicture}

\figcap{电阻的串联与并联}\label{appE:fig:s22-legacy-07}
\end{center}}

\newcommand{\circuitTerminologyDiagram}{%
\begin{center}
\begin{tikzpicture}[x=1cm,y=1cm,>=Latex]
  \coordinate (A) at (-3.6,1.35); \coordinate (B) at (0,1.35);
  \coordinate (C) at (3.6,1.35); \coordinate (D) at (-3.6,-1.35);
  \coordinate (E) at (0,-1.35); \coordinate (F) at (3.6,-1.35);
  \foreach \p/\n in {A/$a$,B/$b$,C/$c$,D/$d$,E/$e$,F/$f$}{
    \fill[DeepBlue] (\p) circle (2.7pt) node[above=4pt] {\n};
  }
  \draw[thick] (A)--node[above] {元件 1}(B)--node[above] {元件 2}(C);
  \draw[thick] (D)--node[below] {元件 4}(E)--node[below] {元件 5}(F);
  \draw[thick] (A)--node[left] {支路}(D);
  \draw[thick] (B)--node[right] {元件 3}(E);
  \draw[thick] (C)--(F);
  \draw[DeepBlue,very thick,rounded corners=3pt,-{Latex[length=2mm]}]
    (-3.95,1.7)--(0.35,1.7)--(0.35,-1.0)--(-3.95,-1.0)--cycle;
  \node[DeepBlue] at (-1.8,0.25) {一个回路};
  \node[TextGray,align=left,anchor=west] at (4.45,0.55)
    {节点：支路的连接点\\支路：相邻节点间的通路\\回路：从节点出发并返回的闭合路径};
\end{tikzpicture}

\figcap{电路网络中的元件、节点、支路与回路}\label{appE:fig:s22-legacy-08}
\end{center}}

\newcommand{\kirchhoffDiagram}{%
\begin{center}
\begin{tikzpicture}[scale=0.9,>=Latex]
  \fill[DeepBlue] (0,0) circle (3pt) node[below right] {节点};
  \draw[vector] (-3.2,1.4)--(-0.2,0.1) node[midway,above] {$I_1$};
  \draw[vector] (-3.2,-1.4)--(-0.2,-0.1) node[midway,below] {$I_2$};
  \draw[vector] (0.2,0.12)--(3.2,1.35) node[midway,above] {$I_3$};
  \draw[vector] (0.2,-0.12)--(3.2,-1.35) node[midway,below] {$I_4$};
  \node[TextGray] at (0,-2.15) {KCL：$I_1+I_2=I_3+I_4$};
  \draw[DeepBlue,thick,rounded corners] (5.1,-1.25) rectangle (9.2,1.25);
  \draw[vector] (5.45,-0.85)--(7.0,-0.85);
  \node[TextGray] at (7.15,-2.15) {KVL：沿闭合回路电压代数和为零};
\end{tikzpicture}

\figcap{基尔霍夫电流定律与电压定律}\label{appE:fig:s22-legacy-09}
\end{center}}

\newcommand{\nodeMethodDiagram}{%
\begin{center}
\begin{tikzpicture}[x=1cm,y=1cm,>=Latex]
  \coordinate (a) at (-3.4,1.35); \coordinate (b) at (0,1.35); \coordinate (c) at (3.4,1.35);
  \coordinate (g1) at (-3.4,-1.35); \coordinate (g2) at (0,-1.35); \coordinate (g3) at (3.4,-1.35);
  \draw[thick] (a)--(b)--(c); \draw[thick] (g1)--(g2)--(g3);
  \draw[thick] (a)--(-3.4,0.55);
  \draw[very thick,DeepBlue] (-3.75,0.38)--(-3.05,0.38) (-3.62,0.03)--(-3.18,0.03);
  \draw[thick] (-3.4,0.03)--(g1); \node[left] at (-3.85,0.2) {$\mathcal E_0$};
  \draw[thick,decorate,decoration={zigzag,segment length=5pt,amplitude=2.2pt}] (b)--(g2);
  \draw[thick,decorate,decoration={zigzag,segment length=5pt,amplitude=2.2pt}] (c)--(g3);
  \draw[thick,decorate,decoration={zigzag,segment length=5pt,amplitude=2.2pt}] (a)--(b);
  \node[above] at (-1.7,1.65) {$R_1$}; \node[right] at (0.35,0) {$R_2$}; \node[right] at (3.75,0) {$R_3$};
  \fill[DeepBlue] (a) circle (2.7pt) node[above left] {$a:\ \varphi_a=\mathcal E_0$};
  \fill[DeepBlue] (b) circle (2.7pt) node[above right] {$b:\ \varphi_b=\varphi$};
  \fill[Ink] (g2) circle (2.7pt) node[below] {参考节点 $\varphi=0$};
  \node[TextGray,align=left,anchor=west] at (5.0,0.45)
    {选参考节点后，各支路电压\\自动等于节点电势之差；\\只需对未知节点列 KCL};
\end{tikzpicture}

\figcap{节点电压法：先求节点电势，再由伏安特性求支路电流}\label{appE:fig:s22-legacy-10}
\end{center}}

\newcommand{\tolmanDiagram}{%
\begin{center}
\begin{tikzpicture}[scale=0.9,>=Latex]
  \draw[curve] (0,0) ellipse (2.25 and 1.35);
  \draw[vector] (0,0)--(0.95,0.8) node[midway,above left] {$r$};
  \draw[DeepBlue,very thick,-{Latex[length=2.5mm]}] (1.65,0.92) arc[start angle=35,end angle=135,x radius=2.05,y radius=1.28];
  \node[DeepBlue] at (0,1.95) {$\omega$};
  \draw[vector] (2.25,0)--(2.25,-1.05) node[right] {$\vect v_e(0)$};
  \draw[very thick] (-0.15,1.35)--(0.15,1.35);
  \node[TextGray,align=left,anchor=west] at (4.0,0.55)
    {晶格突然停止\\电子因惯性继续运动\\随后在碰撞中指数衰减};
\end{tikzpicture}

\figcap{托尔曼--斯图尔特转动铜环实验}\label{appE:fig:s22-legacy-11}
\end{center}}

\newcommand{\transientCurrentDiagram}{%
\begin{center}
\begin{tikzpicture}[scale=0.9,>=Latex]
  \draw[axis] (-0.4,0)--(6.0,0) node[right] {$t$};
  \draw[axis] (0,-0.2)--(0,3.2) node[above] {$I$};
  \draw[guide] (0,2.55)--(5.7,2.55) node[right] {$U/R$};
  \draw[DeepBlue,very thick] (0,0) .. controls (0.75,1.75) and (1.7,2.35) .. (5.5,2.54);
  \draw[<->] (0,-0.55)--node[below] {$\tau$}(1.55,-0.55);
  \node[TextGray] at (2.8,-1.15) {$I(t)=\dfrac UR(1-\e^{-t/\tau})$，载流子惯性使电流不能瞬间建立};
\end{tikzpicture}

\figcap{突加电压后的微观暂态电流}\label{appE:fig:s22-legacy-12}
\end{center}}

\newcommand{\rotatingRingDiagram}{%
\begin{samepage}
\begin{center}
\begin{tikzpicture}[scale=0.70,>=Latex]
  \draw[curve] (0,0) circle (1.65);
  \foreach \a in {0,30,...,330}{\node[DeepBlue] at ({1.65*cos(\a)},{1.65*sin(\a)}) {$+$};}
  \draw[DeepBlue,very thick,-{Latex[length=2.5mm]}] (1.35,0.95) arc[start angle=35,end angle=135,radius=1.65];
  \draw[vector] (0,0)--(1.1,0) node[midway,above] {$R$};
  \node[TextGray,align=left,anchor=west] at (3.4,0.6)
    {每转一周，任一截面通过总电荷 $Q$\\周期 $T=2\pi/\omega$};
\end{tikzpicture}

\figcap{转动带电环与等效环电流}\label{appE:fig:s22-legacy-13}
\end{center}
\end{samepage}}

\newcommand{\vacuumDiodeDiagram}{%
\begin{center}
\begin{tikzpicture}[scale=0.86,>=Latex]
  \draw[DeepBlue,very thick] (-3.5,-1.2)--(-3.5,1.2) node[above] {阴极 $\varphi=0$};
  \draw[DeepBlue,very thick] (3.5,-1.2)--(3.5,1.2) node[above] {阳极 $\varphi=U$};
  \foreach \x in {-2.7,-1.9,...,2.9}{\node[Ink] at (\x,0) {$-$};}
  \draw[vector] (-2.8,-0.75)--(2.7,-0.75) node[midway,below] {电子运动};
  \draw[axis] (-3.5,-1.75)--(4.1,-1.75) node[right] {$x$};
  \draw[<->] (-3.5,-2.25)--node[fill=white,inner sep=1pt] {$L$}(3.5,-2.25);
  \node[TextGray] at (0,-2.95) {$\varphi(x)=U(x/L)^{4/3}$，空间电荷使电势不是线性分布};
\end{tikzpicture}

\figcap{平板真空二极管中的空间电荷限制电流}\label{appE:fig:s22-legacy-14}
\end{center}}

\newcommand{\nonlinearShellDiagram}{%
\begin{samepage}
\begin{center}
\begin{tikzpicture}[scale=0.70,>=Latex]
  \fill[PaleBlue] (0,0) circle (1.85); \fill[white] (0,0) circle (0.82);
  \draw[curve] (0,0) circle (0.82); \draw[curve] (0,0) circle (1.85);
  \foreach \a in {0,45,...,315}{\draw[component] ({0.92*cos(\a)},{0.92*sin(\a)})--({1.7*cos(\a)},{1.7*sin(\a)});}
  \draw[<->] (0,0)--node[above] {$a$}(0.82,0);
  \draw[<->] (0,-0.35)--node[below] {$b$}(1.85,-0.35);
  \node[TextGray,align=left,anchor=west] at (3.5,0.6)
    {$j(r)=I/(4\pi r^2)$\\$j=kE^2$\\故 $E\propto1/r$};
\end{tikzpicture}

\figcap{电导率满足 $\sigma=kE$ 的厚球壳}\label{appE:fig:s22-legacy-15}
\end{center}
\end{samepage}}

\newcommand{\ionChamberDiagram}{%
\begin{center}
\begin{tikzpicture}[scale=0.83,>=Latex]
  \draw[DeepBlue,very thick] (-4.1,1.35)--(4.1,1.35) node[right] {阴极 $-U_0$};
  \draw[DeepBlue,very thick] (-4.1,-1.35)--(4.1,-1.35) node[right] {阳极};
  \draw[guide] (-3.2,-0.35)--(3.1,-0.35) node[right] {粒子径迹};
  \foreach \x in {-2.7,-1.9,...,2.9}{\node[DeepBlue] at (\x,-0.15) {$+$}; \node[Ink] at (\x,-0.55) {$-$};}
  \draw[vector] (-1.4,-0.05)--(-1.4,0.95) node[right] {$v^+$};
  \draw[vector] (0.8,-0.55)--(0.8,-1.12) node[right] {$v^-$};
  \draw[<->] (-4.55,-1.35)--node[left] {$D$}(-4.55,1.35);
  \draw[<->] (3.55,-1.35)--node[right] {$x_0$}(3.55,-0.35);
  \node[TextGray] at (0,-2.15) {电子先到达阳极形成快分量，正离子随后形成慢分量};
\end{tikzpicture}

\figcap{平行板电离室中的电子--离子漂移}\label{appE:fig:s22-legacy-16}
\end{center}}

\newcommand{\ringResistanceDiagram}{%
\begin{center}
\begin{tikzpicture}[scale=0.9,>=Latex]
  \coordinate (RingA) at (-2.7,0);
  \coordinate (RingO) at (0,0);
  \coordinate (RingB) at (2.13,-1.49);
  \draw[curve] (0,0) circle (1.75);
  \fill[DeepBlue] (-1.75,0) circle (3pt) node[left] {$a$};
  \fill[DeepBlue] ({1.75*cos(-35)},{1.75*sin(-35)}) circle (3pt) node[right] {$b$};
  \draw[axis] (0,0)--(-2.7,0);
  \draw[axis] (0,0)--({2.6*cos(-35)},{2.6*sin(-35)});
  \draw pic[draw,"$\theta$",angle radius=11mm,angle eccentricity=1.35]
    {angle=RingA--RingO--RingB};
  \node[TextGray,align=left,anchor=west] at (4.1,0.5)
    {两段电阻 $R_1+R_2=R$\\在端点 $a,b$ 之间并联};
\end{tikzpicture}

\figcap{圆环电阻被接点分成两条并联支路}\label{appE:fig:s22-legacy-17}
\end{center}}

\newcommand{\dualSourceDiagram}{%
\begin{center}
\begin{tikzpicture}[scale=0.88,>=Latex]
  \draw[thick] (-4.0,1.5)--(4.0,1.5) (-4.0,-1.5)--(4.0,-1.5);
  \draw[thick] (-3.2,1.5)--(-3.2,0.95);
  \draw[very thick,DeepBlue] (-3.55,0.75)--(-2.85,0.75) (-3.43,0.38)--(-2.97,0.38);
  \draw[thick] (-3.2,0.38)--(-3.2,-0.1);
  \draw[thick,decorate,decoration={zigzag,segment length=5pt,amplitude=2.2pt}] (-3.2,-0.1)--(-3.2,-1.05);
  \draw[thick] (-3.2,-1.05)--(-3.2,-1.5);
  \node[left] at (-3.65,0.55) {$\mathcal E_1$}; \node[left] at (-3.65,-0.55) {$r_1$};
  \draw[thick] (0,1.5)--(0,0.95);
  \draw[very thick,DeepBlue] (-0.35,0.75)--(0.35,0.75) (-0.23,0.38)--(0.23,0.38);
  \draw[thick] (0,0.38)--(0,-0.1);
  \draw[thick,decorate,decoration={zigzag,segment length=5pt,amplitude=2.2pt}] (0,-0.1)--(0,-1.05);
  \draw[thick] (0,-1.05)--(0,-1.5);
  \node[right] at (0.45,0.55) {$\mathcal E_2$}; \node[right] at (0.45,-0.55) {$r_2$};
  \draw[thick] (3.2,1.5)--(3.2,0.95);
  \draw[thick,decorate,decoration={zigzag,segment length=5pt,amplitude=2.2pt}] (3.2,0.95)--(3.2,-0.95);
  \draw[thick] (3.2,-0.95)--(3.2,-1.5);
  \node[right] at (3.55,0) {$R$};
  \draw[vector] (2.75,1.85)--(3.65,1.85) node[midway,above] {$I$};
\end{tikzpicture}

\figcap{两个有内阻电源并联向负载供电}\label{appE:fig:s22-legacy-18}
\end{center}}


\begin{origexamplebox}{【真空二极管】}\label{appE:ex:s22-33-4}
平板真空二极管的阴极、阳极面积均为 $S$，间距为 $L$。阴极位于 $x=0$ 且电势为零，阳极位于
$x=L$ 且电势为 $U>0$。电子从阴极以可忽略的初速度发射，稳恒电流为 $I$。
\begin{enumerate}[label=(\arabic*),leftmargin=2.2em]
  \item 写出板间电势满足的方程；
  \item 用 $\varphi(x)$ 表示电子速度；
  \item 用 $I$、速度表示电子数密度；
  \item 消去速度和数密度，写出电势的微分方程与空间电荷限制状态的边界条件；
  \item 求伏安关系并证明查尔德--朗缪尔定律 $I=KU^{3/2}$。
\end{enumerate}

\QAENamed{\QAEVacuumQuestion}{阴极、阳极、电势边界与电子运动方向；约定电流大小为正。}{quality:appE:q11}
\end{origexamplebox}

\begin{xesolution}
\textbf{(1)} 电子数密度为 $n(x)$，体电荷密度 $\rho_q=-en$。泊松方程给出
\begin{equation}
  \frac{\dd^2\varphi}{\dd x^2}=-\frac{\rho_q}{\varepsilon_0}
  =\frac{en}{\varepsilon_0}.
\end{equation}

\textbf{(2)} 电子从阴极由静止出发。能量守恒
\[
  \frac12mv^2=e\varphi,
\]
所以
\begin{equation}
  \boxed{v(x)=\sqrt{\frac{2e\varphi(x)}m}}.
\end{equation}

\textbf{(3)} 稳恒电流在各截面相同，取电流大小有
\begin{equation}
  I=en(x)v(x)S,
  \qquad
  \boxed{n(x)=\frac{I}{eSv(x)}}.
\end{equation}

\textbf{(4)} 消去 $n,v$：
\begin{equation}
  \frac{\dd^2\varphi}{\dd x^2}
  =\frac{I}{\varepsilon_0S}\sqrt{\frac{m}{2e}}\,\varphi^{-1/2}.
\end{equation}
空间电荷限制的最大稳定电流状态下，阴极表面电场恰为零；若阴极场反向，发出的电子会被推回。
边界条件为
\begin{equation}
  \varphi(0)=0,
  \qquad \varphi'(0)=0,
  \qquad \varphi(L)=U.
\end{equation}

\textbf{(5)} 记
\[
  A_0=\frac{I}{\varepsilon_0S}\sqrt{\frac{m}{2e}}.
\]
用 $\varphi'$ 乘微分方程并从阴极积分：
\[
  \frac12(\varphi')^2=2A_0\sqrt\varphi,
  \qquad
  \varphi'=2\sqrt{A_0}\,\varphi^{1/4}.
\]
再次积分得
\[
  \varphi(x)=\left(\frac32\sqrt{A_0}\,x\right)^{4/3}.
\]
代入 $\varphi(L)=U$，得到电势分布和电流：
\begin{equation}
  \boxed{\varphi(x)=U\left(\frac xL\right)^{4/3}},
\end{equation}
\begin{equation}
  \boxed{I=\frac49\varepsilon_0S\sqrt{\frac{2e}{m}}
  \frac{U^{3/2}}{L^2}}.
\end{equation}
因此
\begin{equation}
  \boxed{K=\frac49\frac{\varepsilon_0S}{L^2}\sqrt{\frac{2e}{m}}}.
\end{equation}
这一定律建立在平板近似、零初速度、充分电子发射和稳恒一维空间电荷等条件上。
图~\ref{quality:appE:a11} 中电子从零电势阴极向高电势阳极运动；$\varphi\propto x^{4/3}$ 也表明板间电场
$E_x=-\varphi'$、$|\vect E|=\varphi'$ 并非均匀场。

\QAENamed{\vacuumDiodeDiagram}{电子向高电势端加速；电势为非线性分布，电场指向阴极。}{quality:appE:a11}

图~\ref{quality:appE:a11}中电子速率沿 $+x$ 增大，电势也随 $x$ 增大，但 $E_x=-\varphi'<0$，负电荷所受电力仍沿 $+x$。必须区分电场方向、电子速度方向和电流的大小。
\end{xesolution}

\begin{origexamplebox}{【电离计数器】}\label{appE:ex:s22-33-6}
电离计数器测量带电粒子穿过气体时产生的电离电荷。电场把电离产物中的电子和正离子分别引向
阳极和阴极；若入射粒子在气体中被完全吸收，还可由总电离量估计粒子能量。最简单的电离室由
一对平行板电极构成，密封容器中充有工作气体。板间距为 $D$、面积为 $S$，阴极电势为 $-U_0$，
阳极接读出电路。带电粒子在距阳极 $x_0$ 处沿平行于极板的径迹产生 $N$ 对电子--离子；若粒子
沉积能量为 $E_{\rm dep}$、气体中产生一对电子--离子所需平均能量为 $W$，则
$N=E_{\rm dep}/W$。正离子与电子的漂移速率分别为 $v^+,v^-$，假设电荷全部被收集，忽略复合、
电子俘获和次级电离。
\begin{enumerate}[label=(\arabic*),leftmargin=2.2em]
  \item 求均匀电场；
  \item 求极板电容；
  \item 当负载电阻很大时，求输出电压变化 $\Delta U(t)$；
  \item 在 $v^-\gg v^+$ 时描述波形；
  \item 若负载电阻 $R$ 不可忽略，写出输出电压的时间关系。
\end{enumerate}

采用小信号近似，阳极初始电势为零，脉冲幅值 $\Delta U=-\delta V_{\rm anode}$；要求 $|\delta V_{\rm anode}|\ll U_0$，从而可将漂移速率视为常量。

\QAENamed{\QAEIonQuestion}{平行板电离室中电子与离子的两段漂移距离及读出电极。}{quality:appE:q12}
\end{origexamplebox}

\begin{xesolution}
取 $x$ 从阳极指向阴极，阳极 $x=0$、阴极 $x=D$。

\textbf{(1)} 忽略边缘效应，
\begin{equation}
  \boxed{E=\frac{U_0}{D}}.
\end{equation}
场方向由阳极指向阴极；电子向阳极漂移，正离子向阴极漂移。

\textbf{(2)} 气体介电常数近似为 $\varepsilon_0$，
\begin{equation}
  \boxed{C=\frac{\varepsilon_0S}{D}}.
\end{equation}

\textbf{(3)} 平行板的读出权重场为 $1/D$。每个仍在漂移的电荷产生大小为 $ev/D$ 的感应电流。
令
\[
  t_- =\frac{x_0}{v^-},
  \qquad
  t_+ =\frac{D-x_0}{v^+}
\]
分别为电子到阳极、离子到阴极的时间。感应电流大小为
\begin{equation}
  i_{\rm ind}(t)=\frac{Ne}{D}
  \left[v^-H(t_--t)+v^+H(t_+-t)\right],
\end{equation}
其中 $H$ 为阶跃函数。负载近似开路时，感应电荷全部给电容充电：
\begin{equation}
  \boxed{\Delta U(t)=\frac{Ne}{CD}\left[
  v^-\min(t,t_-)+v^+\min(t,t_+)\right]}.
\end{equation}
当两类载流子全部到达电极，$\Delta U(\infty)=Ne/C$，与电离位置 $x_0$ 无关。

\textbf{(4)} 当电离位置不紧靠阴极且 $t_-\ll t_+$ 时，电子在短时间 $t_-$ 内形成陡峭快分量
$Ne x_0/(CD)$；随后正离子在较长时间 $t_+$ 内形成缓慢上升分量
$Ne(D-x_0)/(CD)$。两部分相加仍为 $Ne/C$。

\textbf{(5)} 负载电阻使已形成的电压以时间常数 $\tau_c=RC$ 泄放。电路方程为
\begin{equation}
  \boxed{C\frac{\dd\Delta U}{\dd t}+\frac{\Delta U}{R}=i_{\rm ind}(t)}.
\end{equation}
取 $\Delta U(0)=0$，统一写成卷积形式
\begin{equation}
  \boxed{\Delta U(t)=\frac1C\int_0^t
  i_{\rm ind}(t')\e^{-(t-t')/(RC)}\dd t'}.
\end{equation}
它表明每一小段新产生的感应电压都从产生时刻起指数衰减。若粒子在气体中完全损失能量
$E_p$，且产生一对电子--离子的平均能量为 $W$，还可用 $N=E_p/W$ 反推粒子能量。
图~\ref{quality:appE:a12} 把两段漂移路程直接标成 $x_0$ 与 $D-x_0$：它们分别决定快分量和慢分量的持续时间
及最终幅度。

\QAENamed{\ionChamberDiagram}{电子与正离子的漂移分别形成快、慢脉冲分量；输出取负电压脉冲的正幅值。}{quality:appE:a12}

图~\ref{quality:appE:a12}中电子向阳极、正离子向阴极的漂移对负电压脉冲幅值均作正贡献。只有两类载流子全部到达后，路径长度之和才为 $D$，于是开路总幅值与 $x_0$ 无关；有限负载下较早形成的分量已泄放，峰值通常重新依赖位置。
\end{xesolution}

\begin{origexamplebox}{【圆环电阻】}\label{appE:ex:s22-33-7}
总电阻为 $R$ 的均匀电阻丝弯成圆环，在环上两点 $a,b$ 之间接入电源。固定 $a$，移动 $b$。
\begin{enumerate}[label=(\arabic*),leftmargin=2.2em]
  \item 若电源保持电压 $U$，何时总发热功率最小，最小值是多少；
  \item 若改用保持电流 $I$ 的理想恒流源，结果怎样？
\end{enumerate}

\QAENamed{\QAERingQuestion}{圆环的两个电源接点；接点间不另接短路导线。}{quality:appE:q13}
\end{origexamplebox}

\begin{xesolution}
两接点把圆环分成电阻 $R_1,R_2$ 的两段，且 $R_1+R_2=R$。两段在 $a,b$ 之间并联：
\begin{equation}
  R_{\rm eq}=\frac{R_1R_2}{R_1+R_2}=\frac{R_1R_2}{R}.
\end{equation}

\textbf{(1) 固定电压。}
\[
  P_U=\frac{U^2}{R_{\rm eq}}=\frac{U^2R}{R_1R_2}.
\]
在 $R_1+R_2=R$ 下，乘积 $R_1R_2$ 在 $R_1=R_2=R/2$ 时最大，故功率最小。此时 $b$ 与 $a$
位于直径两端，
\begin{equation}
  \boxed{R_{\rm eq}=\frac R4},
  \qquad
  \boxed{P_{U,\min}=\frac{4U^2}{R}}.
\end{equation}

\textbf{(2) 固定电流。}
\[
  P_I=I^2R_{\rm eq}=I^2\frac{R_1R_2}{R}.
\]
此时要使乘积 $R_1R_2$ 尽量小，因此 $b$ 应尽量靠近 $a$。在允许两接点重合的理想模型中
\begin{equation}
  \boxed{P_{I,\min}=0\quad(b=a)}.
\end{equation}
若规定 $a,b$ 必须不同，则没有被实际取到的正最小值，只有当 $b\to a$ 时趋于零的下确界。固定
电压与固定电流给出相反的优化方向，是因为一个最小化 $1/R_{\rm eq}$，另一个最小化
$R_{\rm eq}$。
图~\ref{quality:appE:a13} 中两条圆弧正是并联的 $R_1,R_2$；把 $b$ 移到 $a$ 的对径点，两条圆弧等长，
对应 $R_1=R_2$。

\QAENamed{\ringResistanceDiagram}{两段圆弧是同端口并联支路；优化变量满足 $R_1+R_2=R$。}{quality:appE:a13}

图~\ref{quality:appE:a13}中的两个支路共用 $a,b$，即使两支路沿不同圆弧也仍是并联。恒流和恒压分别约束电源的不同量，所以同一个 $R_{\rm eq}$ 的极大值在两种优化中具有相反的作用。
\end{xesolution}

\endgroup

\Needspace{18\baselineskip}

\section{复杂电阻网络}

\begingroup

\graphicspath{{../../../sample/xes-original/twenty-third-section/}{../../../sample/xes-original/twenty-third-section/assets/}}

\tikzset{
  resistor/.style={thick,decorate,decoration={zigzag,segment length=4.2pt,amplitude=1.8pt}},
  circuitwire/.style={thick,Ink},
  samenode/.style={circle,draw=DeepBlue,fill=PaleBlue,inner sep=1.6pt},
  terminal/.style={circle,fill=DeepBlue,inner sep=2.2pt}
}

\newcommand{\sourceEquivalentDiagram}{%
\begin{center}
\begin{tikzpicture}[scale=0.88,>=Latex]
  \begin{scope}[xshift=-3.4cm]
    \draw[circuitwire] (-1.5,1)--(-0.45,1);
    \draw[very thick,DeepBlue] (-0.45,1.35)--(-0.45,0.65) (-0.08,1.24)--(-0.08,0.76);
    \draw[resistor] (-0.08,1)--(1.25,1);
    \draw[circuitwire] (1.25,1)--(1.55,1) (-1.5,-1)--(1.55,-1);
    \node at (-0.25,1.65) {$\mathcal E$}; \node at (0.65,1.45) {$r$};
    \node[terminal] at (-1.5,1){}; \node[terminal] at (-1.5,-1){};
  \end{scope}
  \draw[<->,DeepBlue,very thick] (-1.0,0)--(1.0,0)
    node[midway,above] {$\mathcal E=I_{\rm N}r$};
  \begin{scope}[xshift=3.4cm]
    \draw[circuitwire] (-1.5,1)--(1.45,1) (-1.5,-1)--(1.45,-1);
    \draw[circuitwire] (0.15,1)--(0.15,0.45) (0.15,-0.45)--(0.15,-1);
    \draw[resistor] (0.15,0.45)--(0.15,-0.45);
    \draw[circuitwire] (1.05,1)--(1.05,0.48) (1.05,-0.48)--(1.05,-1);
    \draw[DeepBlue,very thick,->] (1.05,-0.48)--(1.05,0.48);
    \node[left] at (-0.05,0) {$r$}; \node[right] at (1.25,0) {$I_{\rm N}$};
    \node[terminal] at (-1.5,1){}; \node[terminal] at (-1.5,-1){};
  \end{scope}
\end{tikzpicture}

\figcap{实际电压源与实际电流源的端口等效}\label{appE:fig:s23-legacy-01}
\end{center}}

\newcommand{\theveninDiagram}{%
\begin{center}
\begin{tikzpicture}[scale=0.88,>=Latex]
  \draw[rounded corners=2mm,fill=PaleBlue,draw=RuleGray] (-4,-1.35) rectangle (-1.25,1.35);
  \node[align=center] at (-2.62,0) {任意线性\\有源二端网络};
  \draw[circuitwire] (-1.25,0.75)--(-0.65,0.75) (-1.25,-0.75)--(-0.65,-0.75);
  \node[terminal] at (-0.65,0.75){}; \node[terminal] at (-0.65,-0.75){};
  \draw[->,DeepBlue,very thick] (0,0)--(1.1,0);
  \draw[circuitwire] (1.6,0.75)--(2.1,0.75);
  \draw[very thick,DeepBlue] (2.1,1.1)--(2.1,0.4) (2.45,0.98)--(2.45,0.52);
  \draw[resistor] (2.45,0.75)--(3.75,0.75);
  \draw[circuitwire] (3.75,0.75)--(4.1,0.75) (1.6,-0.75)--(4.1,-0.75);
  \node at (2.3,1.35) {$U_{\rm oc}$}; \node at (3.15,1.2) {$R_{\rm th}$};
  \node[terminal] at (4.1,0.75){}; \node[terminal] at (4.1,-0.75){};
\end{tikzpicture}

\figcap{戴维南等效只保留开路电压与端口内阻}\label{appE:fig:s23-legacy-02}
\end{center}}

\newcommand{\bridgeDiagram}{%
\begin{center}
\begin{tikzpicture}[scale=0.9,>=Latex]
  \coordinate (L) at (-3,0); \coordinate (T) at (0,1.65);
  \coordinate (B) at (0,-1.65); \coordinate (R) at (3,0);
  \draw[resistor] (L)--node[above left] {$R_1$}(T);
  \draw[resistor] (T)--node[above right] {$R_2$}(R);
  \draw[resistor] (L)--node[below left] {$R_3$}(B);
  \draw[resistor] (B)--node[below right] {$R_4$}(R);
  \draw[resistor,DeepBlue] (T)--node[right] {$R_5$}(B);
  \fill[DeepBlue] (L) circle (2.5pt) node[left] {$A$};
  \fill[DeepBlue] (R) circle (2.5pt) node[right] {$B$};
  \node[TextGray,align=left,anchor=west] at (4.0,0.35)
    {$R_1/R_2=R_3/R_4$\\$\Rightarrow\ U_T=U_B$\\$\Rightarrow\ I_5=0$};
\end{tikzpicture}

\figcap{平衡电桥的中间支路可断开，也可将两端等势点合并}\label{appE:fig:s23-legacy-03}
\end{center}}

\newcommand{\starDeltaDiagram}{%
\begin{center}
\begin{tikzpicture}[scale=0.86,>=Latex]
  \begin{scope}[xshift=-3.1cm]
    \coordinate (A) at (0,1.65); \coordinate (B) at (-1.7,-1.1); \coordinate (C) at (1.7,-1.1);
    \draw[resistor] (A)--node[left] {$R_{AB}$}(B);
    \draw[resistor] (B)--node[below] {$R_{BC}$}(C);
    \draw[resistor] (C)--node[right] {$R_{CA}$}(A);
    \node[above] at (A) {$A$}; \node[left] at (B) {$B$}; \node[right] at (C) {$C$};
  \end{scope}
  \draw[<->,DeepBlue,very thick] (-0.55,0)--(0.55,0);
  \begin{scope}[xshift=3.1cm]
    \coordinate (A) at (0,1.65); \coordinate (B) at (-1.7,-1.1); \coordinate (C) at (1.7,-1.1); \coordinate (O) at (0,0);
    \draw[resistor] (A)--node[right] {$R_A$}(O);
    \draw[resistor] (B)--node[above left] {$R_B$}(O);
    \draw[resistor] (C)--node[above right] {$R_C$}(O);
    \node[above] at (A) {$A$}; \node[left] at (B) {$B$}; \node[right] at (C) {$C$};
    \fill[DeepBlue] (O) circle (2.2pt);
  \end{scope}
\end{tikzpicture}

\figcap{三角形连接与星形连接的端口等效}\label{appE:fig:s23-legacy-04}
\end{center}}

\newcommand{\gridDiagram}{%
\begin{center}
\begin{tikzpicture}[scale=0.8]
  \foreach \x in {0,1,2}{\draw[curve] (\x,0)--(\x,2);}
  \foreach \y in {0,1,2}{\draw[curve] (0,\y)--(2,\y);}
  \foreach \x in {0,1,2}{\foreach \y in {0,1,2}{\fill[Ink] (\x,\y) circle (1.5pt);}}
  \node[left] at (0,0) {$A$}; \node[right] at (2,0) {$B$}; \node[right] at (2,2) {$C$};
  \draw[->,DeepBlue] (3.0,1)--(4.25,1);
  \node[TextGray,align=left,anchor=west] at (4.6,1)
    {先按端点对称性分组\\再合并等势节点};
\end{tikzpicture}

\figcap{田字格网络的节点与端点标记}\label{appE:fig:s23-legacy-05}
\end{center}}

\newcommand{\tetrahedronDiagram}{%
\begin{center}
\begin{tikzpicture}[scale=0.88]
  \coordinate (A) at (0,2); \coordinate (B) at (-2,-1.25);
  \coordinate (C) at (2,-1.25); \coordinate (D) at (0,-0.35);
  \foreach \u/\v in {A/B,A/C,A/D,B/C,B/D,C/D}{\draw[curve] (\u)--(\v);}
  \foreach \u/\pos in {A/above,B/left,C/right,D/below}{\fill[DeepBlue] (\u) circle (2.3pt) node[\pos] {$\u$};}
  \node[TextGray,anchor=west,align=left] at (3.0,0.35)
    {从 $A$ 流入、从 $B$ 流出时\\$C,D$ 关于交换对称，故等势};
\end{tikzpicture}

\figcap{正四面体六条棱均接电阻 $R$}\label{appE:fig:s23-legacy-06}
\end{center}}

\newcommand{\completeGraphDiagram}{%
\begin{center}
\begin{tikzpicture}[scale=0.82]
  \foreach \i in {0,...,4}{\coordinate (P\i) at ({90+72*\i}:1.8);}
  \foreach \i in {0,...,4}{\foreach \j in {0,...,4}{\ifnum\i<\j\draw[curve] (P\i)--(P\j);\fi}}
  \foreach \i in {0,...,4}{\fill[DeepBlue] (P\i) circle (2.2pt);}
  \node[TextGray,anchor=west,align=left] at (3.0,0.35)
    {$K_N$ 中任意两点都有一只 $R$\\其余 $N-2$ 个节点彼此等势};
\end{tikzpicture}

\figcap{$N=5$ 时的完全图电阻网络}\label{appE:fig:s23-legacy-07}
\end{center}}

\newcommand{\octagonDiagram}{%
\begin{center}
\begin{tikzpicture}[scale=0.75]
  \foreach \i in {0,...,7}{\coordinate (O\i) at ({22.5+45*\i}:2.25);}
  \foreach \i in {0,...,7}{\foreach \j in {0,...,7}{
    \pgfmathtruncatemacro{\d}{mod(\j-\i+8,8)}
    \ifnum\i<\j\ifnum\d=4\else\draw[thin,TextGray] (O\i)--(O\j);\fi\fi}}
  \foreach \i in {0,...,7}{\fill[Ink] (O\i) circle (1.7pt);}
  \fill[DeepBlue] (O0) circle (3pt) node[right] {$A$};
  \fill[DeepBlue] (O4) circle (3pt) node[left] {$B$};
  \node[TextGray,anchor=west,align=left] at (3.7,0.3)
    {删去四条最长对角线\\求对径端点 $A,B$ 间电阻};
\end{tikzpicture}

\figcap{变种八边形：除对径点外任意两点均以 $R$ 相连}\label{appE:fig:s23-legacy-08}
\end{center}}

\newcommand{\cubeDiagram}{%
\begin{center}
\begin{tikzpicture}[scale=0.82]
  \coordinate (A) at (-2,-1.3); \coordinate (B) at (1,-1.3);
  \coordinate (C) at (1,1.5); \coordinate (D) at (-2,1.5);
  \coordinate (E) at (-0.8,-0.35); \coordinate (F) at (2.2,-0.35);
  \coordinate (G) at (2.2,2.45); \coordinate (H) at (-0.8,2.45);
  \foreach \u/\v in {A/B,B/C,C/D,D/A,E/F,F/G,G/H,H/E,A/E,B/F,C/G,D/H}{\draw[curve] (\u)--(\v);}
  \draw[guide] (E)--(F) (E)--(H) (A)--(E);
  \foreach \u in {A,B,C,D,E,F,G,H}{\fill[Ink] (\u) circle (1.7pt);}
  \node[left] at (A) {$A$}; \node[below] at (B) {$B$}; \node[right] at (C) {$C$}; \node[above] at (G) {$C_1$};
  \node[TextGray,anchor=west,align=left] at (3.25,0.55)
    {棱相邻：$A,B$\\面对角：$A,C$\\体对角：$A,G$};
\end{tikzpicture}

\figcap{正方体十二条棱均接电阻 $R$}\label{appE:fig:s23-legacy-09}
\end{center}}

\newcommand{\ladderDiagram}{%
\begin{center}
\begin{tikzpicture}[scale=0.78]
  \node[left] at (0,1.2) {$a$}; \node[left] at (0,-1.2) {$b$};
  \foreach \i in {0,1,2,3}{
    \draw[resistor] (1.45*\i,1.2)--(1.45*\i+1.05,1.2);
    \draw[resistor] (1.45*\i,-1.2)--(1.45*\i+1.05,-1.2);
    \draw[resistor] (1.45*\i+1.05,1.2)--(1.45*\i+1.05,-1.2);
    \fill[Ink] (1.45*\i+1.05,1.2) circle (1.4pt);
    \fill[Ink] (1.45*\i+1.05,-1.2) circle (1.4pt);
  }
  \draw[dotted,thick] (5.7,1.2)--(6.45,1.2) (5.7,-1.2)--(6.45,-1.2);
\end{tikzpicture}

\figcap{每一节含上下两个串联电阻和一只横接电阻}\label{appE:fig:s23-legacy-10}
\end{center}}

\newcommand{\threeRailDiagram}{%
\begin{center}
\begin{tikzpicture}[scale=0.72]
  \foreach \row/\lab in {1.6/A,0/B,-1.6/C}{\node[left] at (0,\row) {$\lab$};}
  \foreach \i in {0,1,2,3}{
    \foreach \row in {1.6,0,-1.6}{\draw[resistor] (1.5*\i,\row)--(1.5*\i+1.05,\row);}
    \draw[resistor] (1.5*\i+1.05,1.6)--(1.5*\i+1.05,0);
    \draw[resistor] (1.5*\i+1.05,0)--(1.5*\i+1.05,-1.6);
  }
  \foreach \row in {1.6,0,-1.6}{\draw[dotted,thick] (5.6,\row)--(6.35,\row);}
\end{tikzpicture}

\figcap{三端无限方格网络及端点 $A,B,C$}\label{appE:fig:s23-legacy-11}
\end{center}}

\newcommand{\nestedTriangleDiagram}{%
\begin{center}
\begin{tikzpicture}[scale=0.85]
  \coordinate (A) at (-3,-1.55); \coordinate (B) at (3,-1.55); \coordinate (C) at (0,3.65);
  \draw[curve] (A)--(B)--(C)--cycle;
  \coordinate (L1) at (-1.5,1.05); \coordinate (R1) at (1.5,1.05); \coordinate (M1) at (0,-1.55);
  \draw[curve] (L1)--(R1)--(M1)--cycle;
  \coordinate (L2) at (-0.75,-0.25); \coordinate (R2) at (0.75,-0.25); \coordinate (M2) at (0,1.05);
  \draw[curve] (L2)--(R2)--(M2)--cycle;
  \coordinate (L3) at (-0.375,0.40); \coordinate (R3) at (0.375,0.40); \coordinate (M3) at (0,-0.25);
  \draw[curve] (L3)--(R3)--(M3)--cycle;
  \node[left] at (A) {$A$}; \node[right] at (B) {$B$}; \node[above] at (C) {$C$};
  \node[TextGray,anchor=west,align=left] at (4.0,0.65)
    {边长依次为\\$a,a/2,a/4,\ldots$};
\end{tikzpicture}

\figcap{相邻内接等边三角形在三个中点处与外层相接}\label{appE:fig:s23-legacy-12}
\end{center}}

\newcommand{\sierpinskiDiagram}{%
\begin{center}
\begin{tikzpicture}[scale=0.78]
  \begin{scope}[xshift=-3.4cm]
    \fill[PaleBlue] (-1.7,-1.1)--(1.7,-1.1)--(0,1.84)--cycle;
    \node[below] at (0,-1.45) {第 0 阶};
  \end{scope}
  \begin{scope}
    \fill[PaleBlue] (-1.7,-1.1)--(0,-1.1)--(-0.85,0.37)--cycle;
    \fill[PaleBlue] (0,-1.1)--(1.7,-1.1)--(0.85,0.37)--cycle;
    \fill[PaleBlue] (-0.85,0.37)--(0.85,0.37)--(0,1.84)--cycle;
    \draw[curve] (-1.7,-1.1)--(1.7,-1.1)--(0,1.84)--cycle;
    \node[below] at (0,-1.45) {第 1 阶};
  \end{scope}
  \begin{scope}[xshift=3.4cm]
    \foreach \sx/\sy in {-1.275/-0.735,0/-0.735,-0.6375/0.3675}{
      \fill[PaleBlue] (\sx-0.425,\sy-0.3675)--(\sx+0.425,\sy-0.3675)--(\sx,\sy+0.3675)--cycle;
    }
    \foreach \sx/\sy in {-0.425/-0.735,0.85/-0.735,0.425/0.3675}{
      \fill[PaleBlue] (\sx-0.425,\sy-0.3675)--(\sx+0.425,\sy-0.3675)--(\sx,\sy+0.3675)--cycle;
    }
    \fill[PaleBlue] (-0.425,0.735)--(0.425,0.735)--(0,1.47)--cycle;
    \draw[curve] (-1.7,-1.1)--(1.7,-1.1)--(0,1.84)--cycle;
    \node[below] at (0,-1.45) {第 2 阶};
  \end{scope}
\end{tikzpicture}

\figcap{谢尔宾斯基镂垫的前几阶构造}\label{appE:fig:s23-legacy-13}
\end{center}}


\begin{origexamplebox}{【正方体电阻】}\label{appE:ex:s23-34-6}
正方体框架的十二条棱均接电阻 $R$。求两顶点分别为下列位置关系时的等效电阻：
\begin{enumerate}[label=(\arabic*),leftmargin=2.2em]
  \item 棱相邻；
  \item 面对角线两端；
  \item 体对角线两端。
\end{enumerate}

\QAENamed{\QAECube}{正方体的十二条电阻棱与八个节点；三种待求端口共用该编号。}{quality:appE:q14}
\end{origexamplebox}

\begin{xesolution}
这里以 $1\,\mathrm A$ 电流法统一处理。将流出端取零电势，按保持两个端点不动的对称变换把其余
顶点分组，列每组的 KCL；流入端电势数值就是等效电阻。

\textbf{(1) 棱相邻。}
除端点外的六个顶点按对称性合并后，只需列少量节点方程。解得流入端电势为
$7R/12$，故
\begin{equation}
  \boxed{R_{\rm edge}=\frac{7}{12}R}.
\end{equation}

\textbf{(2) 面对角。}
交换端点 $A,C$ 的垂直平分镜面使电势变为 $V_A-V$；$B,D,F,H$ 均位于该镜面上，故都为 $V_A/2$，而 $E,G$ 分别保留。节点方程给出
\begin{equation}
  \boxed{R_{\rm face}=\frac34R}.
\end{equation}

\textbf{(3) 体对角。}
与流入端相邻的三个顶点等势，与流出端相邻的三个顶点也等势。网络直接化成
\[
  \frac R3+\frac R6+\frac R3,
\]
其中中间两组之间有六只 $R$ 并联。因此
\begin{equation}
  \boxed{R_{\rm body}=\frac56R}.
\end{equation}
三个结果满足“端点在图上的最短距离越远，等效电阻越大”。图~\ref{quality:appE:a14} 用实线画出全部十二条棱；
透视图中的交叉不增加节点。

\QAENamed{\QAECube}{按图中节点写基尔霍夫电流定律；表中的电势依次对应三组端口。}{quality:appE:a14}

图~\ref{quality:appE:a14}的八节点在每种端口选择下均满足相同邻接关系。令 $V_{\rm out}=0$、注入电流为 $I_0$，以 $I_0R$ 归一化电势 $u_i=V_i/(I_0R)$；流入端满足 $3u_A-\sum_{j\sim A}u_j=1$，其余非流出节点满足 $3u_i-\sum_{j\sim i}u_j=0$。例如 $AB$ 端口有 $3u_A-u_B-u_D-u_E=1$、$3u_C-u_B-u_D-u_G=0$。解该七元线性方程组得到：\par\begin{center}\begin{tabular}{c|cccccccc}端口 & $u_A$ & $u_B$ & $u_C$ & $u_D$ & $u_E$ & $u_F$ & $u_G$ & $u_H$\\\hline $AB$ & $7/12$&$0$&$5/24$&$3/8$&$3/8$&$5/24$&$1/4$&$1/3$\\ $AC$ & $3/4$&$3/8$&$0$&$3/8$&$1/2$&$3/8$&$1/4$&$3/8$\\ $AG$ & $5/6$&$1/2$&$1/3$&$1/2$&$1/2$&$1/3$&$0$&$1/3$\end{tabular}\end{center}注入端的 $u_A$ 就是 $R_{\rm eq}/R$；逐行代回节点式，可独立复核三个答案。
\end{xesolution}

\begin{origexamplebox}{【无穷网络】}\label{appE:ex:s23-34-7}
无限梯形网络如图~\ref{quality:appE:q15}，每只电阻均为 $R$。每一节的上、下支路各串联一只 $R$，随后用一只
$R$ 横接两支路。求最左端 $a,b$ 间的等效电阻。

\QAENamed{\QAELadder{0}}{连接完整的二端梯形网络；横接电阻位于每一节两只串联电阻之后。}{quality:appE:q15}
\end{origexamplebox}

\begin{xesolution}
记整个无限网络的等效电阻为 $X$。越过最左端上下两只串联电阻后，后面的网络与原网络完全相同，
等效电阻仍为 $X$；它与该处横接的电阻 $R$ 并联。因此
\begin{equation}
  X=2R+(R\parallel X)=2R+\frac{RX}{R+X}.
\end{equation}
移项得
\[
  X^2-2RX-2R^2=0.
\]
电阻必须为正，取正根
\begin{equation}
  \boxed{X=(1+\sqrt3)R}.
\end{equation}
图~\ref{quality:appE:a15} 中切去第一节后留下的虚线部分，与切去前的整张图在拓扑和每只电阻值上完全相同；
这正是方程中再次出现 $X$ 的原因。

\QAENamed{\QAELadder{1}}{虚线尾部与原网络相同；第一节之后的横接 $R$ 与尾部 $X$ 并联。}{quality:appE:a15}

图~\ref{quality:appE:a15}的虚线框从第一只横接电阻之后开始，故它与该横接电阻是并联关系，端口最前面的两只 $R$ 才是串联。有限节开路截断从 $X_1=3R$ 出发，以 $X_{n+1}=2R+RX_n/(R+X_n)$ 递推，单调下降且有下界 $2R$，所以极限存在并选择正文的正根。
\end{xesolution}

\begin{origexamplebox}{【三端无穷网络】}\label{appE:ex:s23-34-8}
图~\ref{quality:appE:q16} 的三端无限方格网络中，每只电阻均为 $R$。求 $A,C$ 间和 $A,B$ 间的等效电阻。

\QAENamed{\QAEThreeRail{0}}{连接完整的三端无限条带；左端三个端口之间没有直接的竖直电阻。}{quality:appE:q16}
\end{origexamplebox}

\begin{xesolution}
先把三端无限尾部等效成一个三角形电导网络。以 $1/R$ 为电导单位，由上下反射对称性，设尾部
$AB,BC$ 两边的等效电导均为 $p$，$AC$ 边电导为 $q$。在它左边再添一列三只水平电阻和两只
竖直电阻后，网络应与原网络相同。

对新增的三个内部节点作消元，端口等效电导变换为
\begin{equation}
  p'=\frac{p+1}{3p+4},
\end{equation}
\begin{equation}
  q'=\frac{p^2+2pq+2p+3q+1}{(3p+4)(p+2q+2)}.
\end{equation}
自相似要求 $p'=p, q'=q$。取正电导解，可把最终所需电阻化简为
\begin{equation}
  \boxed{R_{AC}=(1+\sqrt5)R},
\end{equation}
\begin{equation}
  \boxed{R_{AB}=\left(1+\frac{\sqrt5+\sqrt{21}}4\right)R}.
\end{equation}
数值上 $R_{AC}\approx3.236R$，$R_{AB}\approx2.705R$。相邻端点的等效电阻较小，符合电流可在
较短路径间分流的直觉。图~\ref{quality:appE:a16} 把尾部的三个端口都保留下来，所以不能像二端梯形网络那样只设
一个未知电阻；这里必须用两个独立参数 $p,q$ 描述三端口。

\QAENamed{\QAEThreeRail{1}\par\QAEThreeReduction}{三端尾部及其等效三角形；$p/R,q/R$ 均为电导，不能误当电阻。}{quality:appE:a16}

图~\ref{quality:appE:a16}保留尾部的三个端口。还可把三端电导矩阵按上下反射分解：反对称电势型 $(1,0,-1)$ 的特征电导为 $\lambda_1=p+2q$，对称型 $(1,-2,1)$ 的为 $\lambda_2=3p$。一节的两根竖直电阻分别给两模增加 $1,3$，三根水平电阻再各串联单位电阻，故 $\lambda_1=(\lambda_1+1)/(\lambda_1+2)$、$\lambda_2=(\lambda_2+3)/(\lambda_2+4)$。取正根可得 $\lambda_1=(\sqrt5-1)/2$、$\lambda_2=(\sqrt{21}-3)/2$，以及 $p=(\sqrt{21}-3)/6$、$q=(3\sqrt5-\sqrt{21})/12$。端口电流分解到两模后，$R_{AC}/R=2/\lambda_1$，$R_{AB}/R=1/(2\lambda_1)+3/(2\lambda_2)$，从而得到上述两个根式。
\end{xesolution}

\begin{origexamplebox}{【分形电阻】}\label{appE:ex:s23-34-10}
边长为 $a$ 的正三角形 $ABC$，每一步都将每个保留三角形的三边二等分，并挖去中央倒立三角形，
无限进行得到谢尔宾斯基镂垫。
\begin{enumerate}[label=(\arabic*),leftmargin=2.2em]
  \item 求最终剩余面积；
  \item 求所有孔洞的总边长，不计最外边界；
  \item 若保留下来的所有边均由单位长度电阻为 $\rho_\ell$ 的金属丝构成，求 $A,B$ 间等效电阻；
  \item 讨论这些结果所反映的分形性质。
\end{enumerate}

\QAENamed{\QAESierpinski}{谢尔宾斯基网络的有限阶构造；各阶均采用相同的单位长度电阻。}{quality:appE:q17}
\end{origexamplebox}

\begin{xesolution}
设原三角形面积为 $S_0=\sqrt3a^2/4$。

\textbf{(1) 面积。}
第 $n$ 步保留 $3^n$ 个边长 $a/2^n$ 的小三角形，故
\begin{equation}
  S_n=S_0\left(\frac34\right)^n
  \longrightarrow\boxed{0}.
\end{equation}

\textbf{(2) 孔洞总边长。}
第 $n\ge1$ 步新挖出 $3^{n-1}$ 个边长 $a/2^n$ 的孔，新增周长
\[
  L_n=3^{n-1}\cdot3\frac{a}{2^n}=a\left(\frac32\right)^n.
\]
因此
\begin{equation}
  \sum_{n=1}^{\infty}L_n
  =a\sum_{n=1}^{\infty}\left(\frac32\right)^n
  =\boxed{+\infty}.
\end{equation}

\textbf{(3) 等效电阻。}
记第 $n$ 阶有限网络两顶点间电阻为 $X_n$。第 $n+1$ 阶由三个第 $n$ 阶副本组成，每个副本的
线性尺度减半，所以若单位长度电阻不变，其全部电阻同时乘 $1/2$。由于三端对称，一个副本若两端
等效电阻为 $X_n$，其等效三角形每边电阻为 $3X_n/2$；缩小后每边为 $3X_n/4$。把三个这样的
三角形按图~\ref{quality:appE:a17} 连接，对中点节点列 KCL，得到新网络的两端电阻为单个三角形边电阻的 $10/9$，
于是
\begin{equation}
  X_{n+1}=\frac{10}{9}\cdot\frac{3X_n}{4}=\frac56X_n.
\end{equation}
第 0 阶只是三边电阻均为 $\rho_\ell a$ 的三角形，
\[
  X_0=\rho_\ell a\parallel2\rho_\ell a=\frac23\rho_\ell a.
\]
所以
\begin{equation}
  X_n=\frac23\rho_\ell a\left(\frac56\right)^n
  \longrightarrow\boxed{0}.
\end{equation}

\textbf{(4) 讨论。}
面积趋零而内部边界总长发散，表明“占据面积”与“边界复杂度”可以朝相反方向演化；网络电阻的递推则由拓扑因子与长度缩放共同决定，本题每步乘 $5/6$，最终等效电阻趋零。图~\ref{quality:appE:a17} 的前三阶
同时显示：每一步都复制三个半尺度副本，这正是三组几何级数与电阻递推的共同来源。

\QAENamed{\QAESierpinski}{三个半尺度副本组成下一阶；相同单位长度电阻意味着每个副本电阻同时缩小一半。}{quality:appE:a17}

图~\ref{quality:appE:a17}每一步都把三个半尺度副本连接成新网络。若小段电阻不随尺度缩短，同样拓扑的递推因子会是 $5/3$；本题固定单位长度电阻，另有长度因子 $1/2$，乘积才为 $5/6$。结论描述这一特定导线模型，不能推广到不同电阻缩放规则的分形材料。
\end{xesolution}

\endgroup

\Needspace{18\baselineskip}
\section{同题异解中的进一步结果}

下列两项与主章题目采用同一模型，但旧版推导给出了主线解答之外仍有价值的推广，
因此单独保留；主章中的基本结论不在这里重复计算。

\subsection{中性导体球的极化率与非均匀场受力}

第20章已经得到：半径为 $R$ 的中性导体球置于局部匀强电场
$\vect E_0$ 中时，球外附加电势等效于球心偶极子，
\[
 \vect p=4\pi\varepsilon_0R^3\vect E_0.
\]
所以导体球的静电极化率为
\[
 \boxed{\alpha=4\pi\varepsilon_0R^3},
 \qquad \vect p=\alpha\vect E_0.
\]
由球面感应电荷密度
$\sigma(\theta)=3\varepsilon_0E_0\cos\theta$ 还可直接检验中性条件：
\[
 Q_+=\int_0^{2\pi}\!\int_0^{\pi/2}
 3\varepsilon_0E_0\cos\theta\,R^2\sin\theta
 \,\dd\theta\dd\phi
 =3\pi\varepsilon_0E_0R^2,
 \qquad Q_-=-Q_+.
\]
净电荷虽然为零，正、负电荷中心的分离仍产生非零偶极矩。

若外场在小球尺度内缓慢变化，把小球从无场处准静态移到场强 $E$ 处，
建立感应偶极矩所需的相互作用能为
\[
 U_{\rm ind}=-\int_0^E\alpha E'\,\dd E'
 =-\frac12\alpha E^2.
\]
因此低阶受力为
\[
 \boxed{\vect F=-\nabla U_{\rm ind}
 =\frac12\alpha\nabla(E^2)}.
\]
它指向场强较大的区域。该近似要求 $R$ 远小于外场的空间变化尺度；
若外场严格均匀，合力为零，只剩内部应力。

\subsection{两等势导体球的有限间距推广}

% CompletenessThermalFields20261024: finite-distance Kelvin images.
设两个等半径导体球的半径为 $R$、球心距为 $L>2R$，二者导通并保持共同电势 $V$；
无穷远取零电势。以 $O_A,O_B$ 表示球心，沿 $O_AO_B$ 的轴线放置像电荷。
定义 $\alpha>0$ 使 $\cosh\alpha=L/(2R)$。图~\ref{appE:complete:two-sphere-images}
标出第一条反演链：$q_0$ 从 $A$ 球心开始，$q_1$ 位于 $B$ 内，再把 $q_1$ 对 $A$ 反演得到 $q_2$；
另有一条从 $B$ 球心开始的镜面对称链。

\begin{qualityfigure}
\centering
\begin{tikzpicture}[>=Latex,font=\small,scale=1.0]
 \draw[thick](-2,0)circle(1.15);\draw[thick](2,0)circle(1.15);
 \draw[dashed](-3.5,0)--(3.5,0);
 \fill[DeepBlue](-2,0)circle(2.2pt);\node[below]at(-2,-.08){$q_0$};
 \fill[TextGray](1.6694,0)circle(2.2pt);\node[above]at(1.60,.12){$q_1$};
 \fill[DeepBlue](-1.6396,0)circle(2.2pt);\node[above]at(-1.60,.12){$q_2$};
 \node[left]at(-3.15,.5){球 $A$};\node[right]at(3.15,.5){球 $B$};
 \node[below]at(-2,-1.19){$O_A$};\node[below]at(2,-1.19){$O_B$};
 \draw[->,DeepBlue](-2,.08)to[bend left=37]node[above]{$q_1=-Rq_0/L$}(1.6694,.08);
 \draw[->,TextGray](1.65,-.08)to[bend left=28]node[below]{$q_2=-Rq_1/(L-d_1)$}(-1.6396,-.08);
 \draw[<->](-2,-1.65)--node[below]{$L$}(2,-1.65);
 \draw[<->](2,.7)--node[above]{$d_1$}(1.6694,.7);
 \draw[->](-2,0)--(-2.8132,-.8132)node[midway,below left]{$R$};
\end{tikzpicture}
\figcap{有限间距双球的一条交错像电荷链。$d_n$ 从该像所在球的球心向另一球量取；
图按 $R/L=1.15/4$ 绘出 $d_1=R^2/L$ 与 $d_2=R^2/(L-d_1)$。
镜面对称的第二条链未重复画出。\label{appE:complete:two-sphere-images}}
\end{qualityfigure}

\textbf{球面反演为什么给出递推。}
一个距球心 $s>R$ 的外部点电荷 $q$，在该球内的 Kelvin 像为
$q'=-Rq/s$，像点距球心 $d'=R^2/s$；在球面上，两者的电势逐点相消。
从球 $A$ 中心的 $q_0=4\pi\varepsilon_0RV$ 开始，它在 $A$ 球面单独给出常数 $V$。
对 $B$ 反演产生 $q_1$，消除 $q_0$ 在 $B$ 球面的电势；再对 $A$ 反演产生 $q_2$，消除 $q_1$ 在 $A$ 球面的电势。
以后继续交替。若 $d_n$ 是第 $n$ 个像距所在球心、向另一球的距离，则
\begin{equation}
 d_0=0,\qquad
 d_{n+1}=\frac{R^2}{L-d_n},\qquad
 q_{n+1}=-\frac{R}{L-d_n}q_n.
 \label{appE:eq:complete-image-recurrence}
\end{equation}
图中 $q_0$ 与 $q_1$ 在球 $B$ 上相消，$q_1$ 与 $q_2$ 在球 $A$ 上相消，清楚显示了配对方式。
单独这条链给 $A$ 球电势 $V$、$B$ 球电势零；再叠加从 $B$ 球心开始的对称链，才得到两球同时为 $V$ 的边界条件。

\textbf{解出位置和电量。}
利用 $2\cosh\alpha\,\sinh[(n+1)\alpha]-\sinh(n\alpha)=\sinh[(n+2)\alpha]$，
从 $n=0$ 归纳可得
\[
 \frac{d_n}{R}=\frac{\sinh(n\alpha)}{\sinh[(n+1)\alpha]},\qquad
 \frac{q_n}{q_0}=(-1)^n\frac{\sinh\alpha}{\sinh[(n+1)\alpha]}.
\]
例如把第一个位置式代入 $L-d_n$，恰得
$R\sinh[(n+2)\alpha]/\sinh[(n+1)\alpha]$，故同时恢复式~\eqref{appE:eq:complete-image-recurrence}的两个递推。
$L>2R$ 时 $\alpha>0$，$|q_n|$ 按指数衰减；像点留在各球内部，并向
$d_\infty=Re^{-\alpha}<R$ 聚集，所以球外电势级数收敛。
配对相消验证了真实球面的边界条件，无穷远电势也为零；由外部 Dirichlet 问题的唯一性，
该像系统在球外等效于真实导体。包住两球的远处高斯面又说明真实总电荷等于全部像电荷之和。

两条链的总和为 $Q_{\rm tot}=2\sum_{n=0}^\infty q_n$，因此
\[
 \boxed{C(L)=\frac{Q_{\rm tot}}V
 =8\pi\varepsilon_0R\sinh\alpha
 \sum_{n=1}^\infty\frac{(-1)^{n+1}}{\sinh(n\alpha)}}.
\]
电容的量纲为 $\varepsilon_0R$。当 $L\to\infty$，$\alpha\to\infty$，除首项外的比值均趋零，
故 $C(L)\to8\pi\varepsilon_0R$，等于两个远离的孤立球电容之和。

相切极限 $L\to2R^+$ 对应 $\alpha\to0^+$。对每个固定 $n$，
$\sinh\alpha/\sinh(n\alpha)\to1/n$。
这里还要允许极限与级数交换：这些正系数随 $n$ 单调下降，且
$\sinh(n\alpha)\ge n\sinh\alpha$，所以第 $N$ 项之后的交错余项不超过 $1/(N+1)$，
这个界与 $\alpha$ 无关。先控制余项、再对有限项取极限，便有
\[
 C(2R)=8\pi\varepsilon_0R\sum_{n=1}^\infty\frac{(-1)^{n+1}}n
 =\boxed{8\pi\varepsilon_0R\ln2}.
\]
这既恢复第20章的相切结果，也说明有限间距的双曲函数递推怎样连续过渡为 $q_n\propto(-1)^n/(n+1)$。


\endgroup
